\documentclass[12pt]{article}
\PassOptionsToPackage{hidelinks}{hyperref}
\usepackage{pdfpages}
\usepackage{wrapfig}
\usepackage{bm}
\usepackage{upgreek}
\usepackage{dsfont}
\usepackage{simplewick}
\usepackage{mathtools}
\usepackage{framed}
\usepackage{microtype}
\usepackage{marvosym}
\usepackage{color}
\usepackage{transparent}
\usepackage{placeins}
\usepackage{mdframed,mdwlist}
\usepackage{graphicx}
\usepackage{float}
\usepackage{longtable}
\usepackage{mathdots}
\usepackage{eucal}
\usepackage{array}
\usepackage{stmaryrd}
\usepackage{amsthm}
\usepackage{pifont}
\usepackage{lipsum}
\usepackage{enumerate}
\usepackage{amsmath}
\usepackage{amssymb}
\usepackage{algorithm}
\usepackage{algpseudocode}
\usepackage{setspace}
\usepackage{adjustbox}
\usepackage{caption}
\usepackage{orcidlink}
\usepackage[authoryear]{natbib}
\addtolength{\oddsidemargin}{-.5in}%
\addtolength{\evensidemargin}{-1in}%
\addtolength{\textwidth}{1in}%
\addtolength{\textheight}{1.7in}%
\addtolength{\topmargin}{-1in}%


\interfootnotelinepenalty=10000

\DeclareMathOperator*{\argmax}{arg\,max}
\DeclareMathOperator*{\argmin}{arg\,min}
\DeclareMathOperator*{\plim}{plim\,}
\newcommand{\blue}[1]{\textcolor[rgb]{0.00,0.00,1.00}{#1}}
\newenvironment{bluechange}{\begingroup\color[rgb]{0.00,0.00,1.00}}{\endgroup}
\providecommand{\Tilde}[1]{\widetilde{#1}}
\newcommand{\safeincludegraphics}[2][]{%
  \IfFileExists{#2}{\includegraphics[#1]{#2}}{%
    \fbox{\begin{minipage}[c][0.16\textheight][c]{0.8\linewidth}\centering
    Figure file unavailable in the attached source.
    \end{minipage}}%
  }%
}


\newtheorem{theorem}{Theorem}
\newtheorem{proposition}{Proposition}
\newtheorem{lemma}{Lemma}
\newtheorem{corollary}{Corollary}
\theoremstyle{remark}
\newtheorem{remark}{Remark}
\theoremstyle{plain}


\begin{document}
\title{Consistent information criteria for regularised generalised linear models}
\author{ Qingyuan Zhang$^{1\,\orcidlink{0000-0002-0906-8533}}$, Hien Duy Nguyen$^{2,3\,\orcidlink{0000-0002-9958-432X}}$.}
\date{}

\maketitle

\noindent\footnotesize{$^{1}$School of Mathematics and Physics, University of Queensland, St Lucia, QLD 4072, Australia\\
$^{2}$School of Computing, Engineering and Mathematical Sciences, La Trobe University, Bundoora, VIC 3086, Australia\\
$^{3}$Institute of Mathematics for Industry, Kyushu University, Nishi Ward, Fukuoka 819-0395, Japan\\
Email: h.nguyen5@latrobe.edu.au}

\doublespacing
\section*{Abstract}
\begin{bluechange}
We specialise the PanIC framework for information-criterion selection in constrained generalised linear models (GLMs) and give verifiable sufficient conditions for consistent selection of the smallest risk-optimal constraint level. We establish consistency for deterministic continuum-indexed nested families and a transfer result for possibly data-dependent finite candidate sets under an explicit one-sided coverage condition. For regularised linear least squares, we give a Bayesian information criterion (BIC)-like construction and fixed-threshold second-order limit laws, including non-Gaussian boundary limits, together with a subsampling approximation for scalar projections of those limits. We also give a data-dependent calibration rule for the PanIC multiplier and demonstrate the procedures through simulations.
\end{bluechange}

\noindent \textbf{Keywords}: Information criteria; model selection; generalised linear models; least absolute shrinkage and selection operator; regularisation; asymptotic theory



\section{Introduction}
Information criteria balance empirical fit against model complexity. Classical examples include AIC \citep{akaike_new_1974} and BIC \citep{schwarz_estimating_1978}; see \citet{mcquarrie_regression_1998} and \citet{buhlmann_statistics_2011} for broader accounts.

Our study focuses on model selection schemes with consistency guarantees. The work of \cite{sin_information_1996} provides consistency results for a broad class of ICs for generic loss-based learning problems, but its theorems impose conditions on the objective and candidate models that can be difficult to verify. In \cite{nguyen_panic:_2023}, those criteria are studied further for generic risk-minimisation problems, yielding the model-selection framework that we call PanIC. The name uses the Greek prefix ``pan'', meaning ``all'' or ``of everything''. Like \cite{baudry2015estimation}, the PanIC and \cite{sin_information_1996} results allow generic loss functions, whereas the following works use a negative log-likelihood \citep{leroux_consistent_1992, hui_order_2015} or a negative composite, pseudo, or quasi log-likelihood \citep{varin_note_2005, gao_composite_2010, ng_model_2014, hui_use_2021}.

\blue{We specialise the PanIC framework to constrained regression. Consistency of an empirical-risk minimiser within a fixed regression model does not determine the appropriate constraint level. Across a nested constraint family, empirical risk alone favours the largest candidate, which motivates information-criterion selection of the constraint level. We consider Ridge, Lasso, and Elastic-net constraints \citep{hoerl_ridge_1970,tibshirani_regression_1996,zou_regularization_2005}.}

\begin{bluechange}
Recent works have also considered structured high-dimensional regression in which structural or target parameters and nuisance or incidental parameters are regularised differently. For instance, \citet{drikvandi2025} develops smooth and nonsmooth shrinkage procedures for settings with specified or unspecified parameters of interest. Relatedly, \citet{fan_tang_shi_2018} studies a finite-dimensional vector of structural parameters together with a high-dimensional vector of sparse incidental parameters, obtaining consistency and asymptotic normality for the structural-parameter estimators but only partial selection consistency for the incidental parameters. These objectives and asymptotic regimes differ from our fixed-dimensional target of selecting a constraint level in a nested GLM family.

To address the consistency of regularised regression problems, we extend the PanIC and BIC-like information-criteria framework from the finite-model setting of \cite{nguyen_panic:_2023} (see also \citealp{pmlr-v235-westerhout24a} and \citealp{HoNguyen2025VariableSelectionThresholding}) to deterministic continuum-indexed nested model families. A summary of our contributions is as follows:
\begin{enumerate}
    \item We verify sufficient conditions for PanIC consistency in linear, logistic, Poisson, and Gamma regression. For Gamma regression, we also allow the shape parameter to be estimated jointly on a fixed compact interval bounded away from zero.
    \item We establish first-order consistency for deterministic continuum-indexed nested families. We further show that a finite random candidate set, including one constructed from the data, inherits threshold consistency when it contains a candidate approaching the population-optimal threshold from above and the penalty has the required fixed-separation rate.
    \item We give a finite-sample implementation and a repeated-split cross-signed radius rule for the PanIC multiplier that does not use the true coefficient support. A pointwise fixed-dimensional result shows that, whenever the two ordinary half-sample estimators are root-sample-size consistent, the cross-signed $L_1$ target has the same rate. Uniformity over a fixed positive compact multiplier range ensures that calibration does not alter the first-order consistency target. A numerical study evaluates the rule in independent and correlated Gaussian and logistic designs and in an independent Poisson design.
    \item For constrained linear least squares, we derive fixed-threshold second-order laws and the corresponding result after exact selection on a finite grid, including non-Gaussian boundary limits. We also establish a subsampling approximation for nondegenerate scalar projections of the fixed-threshold laws.
\end{enumerate}
\end{bluechange}
Our consistency target is the smallest constraint level attaining the minimum population risk in a nested candidate family. It is distinct from parameter-estimation and support-selection consistency as studied, for example, by \citet{fan_2001} and \citet{zhao_model_2006}. The oracle inequalities of \citet{massart_2011_lasso_l1} also concern a different target: our results establish information-criterion consistency for the stated fixed-dimensional constrained regression families without first deriving finite-sample oracle inequalities.

The rest of this article proceeds as follows: in Section \ref{Sect: Method}, the PanIC and BIC-like IC framework is reviewed and our proposed estimation method is given with theoretical justification. In Section \ref{Sect: Asymptotic-distribution}, we derive second-order asymptotic distribution results for constrained estimators and discuss the non-Gaussian limits that can arise under binding constraints. In Section \ref{Sect: Simulation}, we evaluate the implementable calibration protocol in independent and correlated Gaussian and logistic designs and in a Poisson design, and give a numerical boundary and subsampling illustration. Section \ref{Sect: Conclusion} concludes the paper. Proofs, supporting technical arguments, and reproducibility details are collected in the Appendix.

\section{Method} \label{Sect: Method}
\subsection{PanIC}\label{Sect: PanIC}
Consider the following setting of a general model selection problem. Let $(\Omega, \mathcal{F}, \mathbb{P})$ be the underlying probability space and $X: \Omega \to \mathbb{X} \subset \mathbb{R}^d$ be a random vector with law $\mathbb{P}_X$ on $(\mathbb{X}, \mathcal{B}(\mathbb{X}))$, where $\mathcal{B}(\mathbb{X})$ denotes the Borel $\sigma$-algebra of $\mathbb{X}$. We observe a sequence of independent and identically distributed (i.i.d.) random variables $(X_i)_{i \in \left[n\right]}$ with common distribution $\mathbb{P}_X$, where $\left[n\right] = \{1, \hdots, n\}$. The set of candidate models is given by a sequence of hypotheses $(\mathcal{H}_k)_{k \in \left[m\right]}$. Each $\mathcal{H}_k$ is a parameterised model class indexed by some parameter vector $\beta_k \in \mathbb{T}_k \subset \mathbb{R}^{q_k}$, where $(q_k)_{k \in \left[m\right]} \subset \mathbb{N}$. That is, for each $k \in \left[m\right]$,
\[\mathcal{H}_k = \{h_k(\cdot; \beta_k): \mathbb{X} \to \mathbb{R}, \beta_k \in \mathbb{T}_k\}.\]

Both the functional form of $h_k$ and the parameter space $\mathbb T_k$ may vary with $k\in[m]$. More generally, we evaluate each hypothesis through a loss function
\[
\ell_k:\mathbb{X}\times\mathbb{T}_k\to\mathbb{R},
\]
which may depend on the full observation and on the parameter. For each $k \in \left[m\right]$, define
\[R_{k,n}(\beta_k) = \frac{1}{n} \sum_{i=1}^n \ell_k(X_i; \beta_k)\]
and
\[r_k(\beta_k) = \mathbb{E} \{\ell_k(X; \beta_k)\}\]
as the empirical risk and the expected risk, respectively. This formulation includes the common special case in which $\ell_k(X;\beta_k)=\ell(h_k(X;\beta_k))$ for some scalar loss $\ell$.

Further, we call $k^*$ the optimal hypothesis index,
\begin{equation}\label{Eqs: Optimal hypothesis index}
    k^* = \min\left\{\argmin_{k \in \left[m\right]} \left\{\min_{\beta_k \in \mathbb{T}_k}r_k(\beta_k)\right\}\right\}.
\end{equation}
The corresponding hypothesis space $\mathcal{H}_{k^*}$ is the smallest-indexed candidate class attaining the minimum population risk. Let $\hat{K}_n$ denote an estimator of $k^*$. A model selection scheme is said to be consistent if
\begin{equation} \label{Eqs: Consistency}
    \plim_{n\to\infty} \hat{K}_n = k^*,
\end{equation}
where $\plim$ denotes convergence in probability (as per \citealp[Sec. 3.2]{amemiya1985advanced}). In the case where the set of candidate models is finite, this is equivalent to
\begin{equation}\label{Eqs: Probability Limit}
    \lim_{n\to\infty}\mathbb{P}\left(\hat{K}_n = k^*\right) = 1.
\end{equation}

Now, suppose that the model selection problem satisfies conditions A1--A3, for each $k \in \left[m\right]$:
\begin{enumerate}
    \item[A1]: $\ell_k(x; \beta_k)$ is Carath\'{e}odory in the sense that $\ell_k(x; \cdot): \mathbb{T}_k \to \mathbb{R}$ is continuous for each $x\in \mathbb{X}$, and $\ell_k(\cdot; \beta_k): \mathbb{X} \to \mathbb{R}$ is $\mathcal{B}(\mathbb{X})$-measurable for each $\beta_k\in\mathbb{T}_k$.
    \item[A2]: $\mathbb{T}_k$ is compact and there exists some $\beta_{k,0}\in\mathbb{T}_k$ such that $\ell_k(X; \beta_{k,0})$ is square integrable. That is,
    \[\mathbb{E}\{\ell_k(X;\beta_{k,0})^2\} < \infty.\]
    \item[A3]: There exists some measurable function $\mathcal{G}_k:\mathbb{X}\to \mathbb{R}_{\geq 0}$, such that $\mathbb{E}\{\mathcal{G}_k(X)^2\} < \infty$ and
    \[
    |\ell_k(x; \beta_k) - \ell_k(x;\beta_k')| \leq \mathcal{G}_k(x) \lVert \beta_k - \beta_k'\rVert
    \]
    for all $\beta_k,\beta_k' \in \mathbb{T}_k$ and for almost all $x\in\mathbb{X}$.
\end{enumerate}

A PanIC estimator $\hat{K}_n$ is defined as
\begin{equation}\label{Eqs: PanIC estimator}
    \hat{K}_n = \min\left\{\argmin_{k \in \left[m\right]} \left\{\min_{\beta_k \in \mathbb{T}_k}R_{k,n}(\beta_k) + P_{k,n}\right\}\right\},
\end{equation}
where $P_{k,n}: \Omega \to \mathbb{R}_{>0}$ is allowed to be a stochastic, strictly positive function that satisfies conditions B1 and B2, for each $k \in \left[m\right]$:
\begin{enumerate}
    \item[B1]: $P_{k, n} > 0$ almost surely, for all $n \in \mathbb{N}$, and $P_{k, n}= o_{\mathbb{P}}(1)$ as $n\to\infty$.
    \item[B2]: If $k < l$, then $\plim_{n\to\infty}\sqrt{n}\{P_{l,n}-P_{k,n}\}=\infty$.
\end{enumerate}
Under A1--A3, B1, and B2, Theorem 1 of \cite{nguyen_panic:_2023} establishes consistency of the estimator in (\ref{Eqs: PanIC estimator}).

\cite{nguyen_panic:_2023} also introduced the notion of a BIC-like criterion, one that satisfies the condition $\mathrm{B2}^*$ instead of B2,
\begin{enumerate}
    \item[$\mathrm{B2}^*$]: If $k < l$, then $\plim_{n\to\infty}n\{P_{l,n}-P_{k,n}\}=\infty$.
\end{enumerate}
Theorem 2 of \cite{nguyen_panic:_2023} shows the consistency of BIC-like criteria under additional assumptions C1--C5. For $k \in \left[m\right]$:
\begin{enumerate}
    \item[C1]: $r_k$ is Lipschitz continuous on $\mathbb{T}_k$, and it is twice differentiable and uniquely minimised at some $\beta^*_k\in\mathbb{T}_k$.
    \item[C2]: $\ell_k(x; \cdot)$ is Lipschitz continuous and differentiable at $\beta^*_k$, for almost all $x\in\mathbb{X}$.
    \item[C3]: The set $\mathbb{T}_k$ is second order regular at $\beta^*_k$.
    \item[C4]: The quadratic growth condition holds at $\beta^*_k$.
    \item[C5]: If $r_k(\beta^*_k) = r_{k^*}(\beta^*_{k^*})$, then $n(R_{k,n}(\beta^*_k) - R_{k^*,n}(\beta^*_{k^*})) = O_{\mathbb{P}}(1)$.
\end{enumerate}

Conditions B1, B2, and $\mathrm{B2}^*$ remain satisfied when the penalty function $P_{k,n}$ is multiplied by any fixed positive constant. We will express IC penalties in the form
\begin{equation}\label{Eqs: Penalty Shape}
    P_{k, n} = \kappa \times \mathrm{pen}_{\mathrm{shape}}(k,n)
\end{equation}
and call $\mathrm{pen}_{\mathrm{shape}}:\mathbb{N}\times\mathbb{N}\rightarrow \mathbb{R}_{\geq 0}$ the penalty shape of $P_{k,n}$. The asymptotic conditions therefore leave the fixed multiplier $\kappa$ undetermined; Section \ref{Sect: calibration} gives a prespecified finite-sample calibration rule.

\begin{bluechange}
\begin{proposition}\label{Prop: uniform kappa}
Suppose $m<\infty$ is fixed and a positive penalty shape $Q_{k,n}$ satisfies
\[
\max_{1\le k\le m}Q_{k,n}=o_{\mathbb P}(1),
\qquad
\sqrt n\{Q_{l,n}-Q_{k,n}\}\xlongrightarrow{\mathbb P}\infty
\quad\text{for every }l>k.
\]
Let $0<\underline\kappa\le\overline\kappa<\infty$, and let $\widehat\kappa_n>0$ almost surely be a possibly data-dependent random variable such that
\[
\mathbb P\{\widehat\kappa_n\in[\underline\kappa,\overline\kappa]\}\longrightarrow1.
\]
Then $P_{k,n}=\widehat\kappa_nQ_{k,n}$ satisfies B1 and B2. Consequently, whenever A1--A3 hold, the associated PanIC selector is consistent. After replacing $\sqrt n$ by $n$ in the displayed separation condition, the same argument verifies B1 and $\mathrm{B2}^*$; whenever C1--C5 also hold, the associated BIC-like selector is consistent.
\end{proposition}

Appendix \ref{App: kappa calibration} gives a data-dependent calibration rule whose multiplier remains in the compact range required by the proposition.
\end{bluechange}


\subsection{Regression Problems} \label{Sect: Regression}
In the regression setting, one observation is $Z=(X,Y)\in\mathcal X\times\mathcal Y$. A GLM mean function $h(\cdot;\beta_0,\beta):\mathcal X\to\mathcal M$ is indexed by an intercept $\beta_0\in\mathbb R$ and slopes $\beta\in\mathbb R^d$, and the loss is written $\ell(Z;\beta_0,\beta)$. Table \ref{Table: Regression} gives the linear, logistic, Poisson, and Gamma cases considered below.
\begin{table}[H]
\captionsetup{font=footnotesize}
\caption{Mean and loss functions for generalised linear regression models. Here, $h=h(x;\beta_0,\beta)$. The Gamma model uses a log link, rather than the canonical inverse link.}
\label{Table: Regression}
\centering
\footnotesize
\begin{tabular}{c c c}
 Regression & Mean function & Loss function\\
  [0.5ex]
  \hline\hline
 Linear & $h = \beta_0 + \beta^\top x$ & $\ell(x, y; \beta_0, \beta) = \left(y - h\right)^2$ \rule{0pt}{5ex}\\
 Logistic & $h = (1 + \exp(-(\beta_0 + \beta^\top x)))^{-1}$ & $ \ell(x, y; \beta_0, \beta) = -\left\{y\log h
     + (1 - y)\log\left(1 - h\right)\right\}
$ \rule{0pt}{5ex}\\
 Poisson & $h = \exp(\beta_0 + \beta^\top x)$ & $
 \ell(x, y; \beta_0, \beta) = -\left\{y\log h - h
 -\log(y!) \right\} $ \rule{0pt}{5ex}\\
 Gamma & $h = \exp(\beta_0 + \beta^\top x)$ & $
 \ell(x, y; \beta_0, \beta) = -\left\{ -\nu\log(h)
  + (\nu - 1) \log (y)
 - y\nu/h\right\} $ \rule{0pt}{5ex}\\[4ex]
 \hline\hline
\end{tabular}
\end{table}
\noindent Following GLM practice, we use losses given by the negative log-likelihood up to additive terms independent of $(\beta_0,\beta)$ and, in the Gaussian case, up to a positive multiplicative constant. Additive candidate-independent terms do not affect either the constrained optimiser or the IC comparison. A positive rescaling of the loss preserves each constrained optimiser, but it preserves the IC comparison only when the candidate-dependent penalty is rescaled by the same factor. For the Gamma regression, we use the mean parameterisation of the Gamma density function
\[
p(y;\mu, \nu) = \frac{1}{\Gamma(\nu)}\left(\frac{\nu}{\mu}\right)^\nu y^{\nu-1}\exp\left(-y\frac{\nu}{\mu}\right),
\]
where $\mathbb{E}(Y) = \mu$ and $\mathrm{Var}(Y) = \mu^2/\nu$.
For the baseline formulation, $\nu>0$ may be treated as a fixed shape parameter. Proposition \ref{Prop: estimated Gamma shape} below shows that the PanIC verification also permits joint estimation of $\nu$ when it is restricted to a fixed compact interval bounded away from zero.

Under the Empirical Risk Minimisation principle (ERM), a regression problem is estimated by minimising the empirical risk over the parameter space. Assuming a unique solution exists, learning according to the ERM principle yields the estimates $\hat{\beta}_0$ and $\hat{\beta}$, where
\[
\left(\hat{\beta}_{0},\hat{\beta}\right)=\underset{\beta_{0}\in\mathbb{R},\beta\in\mathbb{R}^{d}}{\arg\min}~\frac{1}{n}\sum_{i=1}^{n}\ell\left(x_{i},y_{i};\beta_{0},\beta\right).
\]
To mitigate the tendency of ERM to overfit, the Lasso regression, the Ridge regression, or the Elastic-net regression are often solved instead. Lasso and Ridge can both be written as follows. For some $C\in\mathbb{R}_{\geq0}$, define
\begin{align}\label{Eqs: Lasso Ridge}
\begin{split}
    \mathcal{H} &= \{h(\cdot; \beta_0, \beta): \mathcal{X} \to \mathcal{M}, (\beta_0, \beta) \in \mathbb{T}\},\\
    \mathbb{T} &= \{(\beta_0, \beta) \, |\, \beta_0 \in \mathbb{U},\, \beta\in \mathbb{R}^d: \lVert\beta\rVert_p^p \leq C\},
\end{split}
\end{align}
where $\mathbb{U}$ is a fixed compact \blue{interval} in $\mathbb{R}$, common to all candidates. \blue{Unless stated otherwise, the population risks and targets are defined relative to this restriction. When the results are interpreted for an otherwise unrestricted-intercept GLM, we assume that every relevant unrestricted population minimiser exists and has its intercept in $\operatorname{int}(\mathbb U)$, so the restriction is locally inactive.} Setting $p=1$ in (\ref{Eqs: Lasso Ridge}) corresponds to the Lasso problem \citep{tibshirani_regression_1996}, and setting $p=2$ corresponds to the Ridge regression problem \citep{hoerl_ridge_1970}.
Slightly modifying (\ref{Eqs: Lasso Ridge}), we also retrieve the Elastic-net problem
\citep{zou_regularization_2005}:
\begin{align}\label{Eqs: Elastic Net}
\begin{split}
    \mathcal{H} &= \{h(\cdot; \beta_0, \beta): \mathcal{X} \to \mathcal{M}, (\beta_0, \beta) \in \mathbb{T}\},\\
    \mathbb{T} &= \{(\beta_0, \beta) \, |\, \beta_0 \in \mathbb{U},\, \beta\in \mathbb{R}^d: \alpha\lVert\beta\rVert_1 + (1-\alpha)\lVert\beta\rVert_2^2 \leq C\},
\end{split}
\end{align}
for some $\alpha \in (0, 1)$ and $C\in\mathbb{R}_{\geq0}$. For each fixed $C$, each of (\ref{Eqs: Lasso Ridge}) and (\ref{Eqs: Elastic Net}) defines one PanIC candidate model. A strictly increasing sequence $(C_k)_{k\in\left[m\right]}$ therefore defines nested parameter spaces: $\mathbb{T}_k\subset\mathbb{T}_l$ whenever $k<l$.
\begin{align}\label{Eqs: Model selection}
\begin{split}
    \mathcal{H}_k &= \{h(\cdot; \beta_0, \beta): \mathcal{X} \to \mathcal{M}, (\beta_0, \beta) \in \mathbb{T}_k\}\\
    \mathbb{T}_k &= \begin{cases}
        \{(\beta_0, \beta) \, |\, \beta_0 \in \mathbb{U}, \beta\in \mathbb{R}^d: \lVert\beta\rVert_1 \leq C_k\}, & \text{for Lasso regression,}\\
        \{(\beta_0, \beta) \, |\, \beta_0 \in \mathbb{U}, \beta\in \mathbb{R}^d: \lVert\beta\rVert_2^2 \leq C_k\}, & \text{for Ridge regression,}\\
        \{(\beta_0, \beta) \, |\, \beta_0 \in \mathbb{U}, \beta\in \mathbb{R}^d: \alpha\lVert\beta\rVert_1 + (1-\alpha)\lVert\beta\rVert_2^2 \leq C_k\}, & \text{for Elastic-net regression}.
    \end{cases}
\end{split}
\end{align}
The PanIC theory can therefore be applied to consistent selection among these nested constraint levels once its loss, compactness, and penalty conditions are verified.

\subsection{Consistency} \label{Sect: Theory}
The following conditions verify A1--A3 for the four regression losses; the Appendix gives the derivations. The consistency conclusion then follows for any penalty satisfying B1 and B2.
\begin{proposition}\label{Prop: Sufficient conditions}
Assume that $(X_i,Y_i)_{i\in\mathbb N}$ is an i.i.d. sequence and that the stated support and moment conditions hold for the regression problem under consideration. Then any PanIC estimator based on a penalty satisfying B1 and B2 enjoys the consistency property in (\ref{Eqs: Consistency}).
    \begin{enumerate}
        \item[1.] Linear regression: $Y\in\mathbb R$, and both $X$ and $Y$ have finite fourth moments.
        \item[2.] Logistic regression: $Y\in\{0,1\}$ almost surely, and $X$ has a finite second moment.
        \item[3.] Poisson regression: $Y\in\mathbb Z_{\geq0}$ almost surely, $X$ is Gaussian or sub-Gaussian, and $Y$ has a finite fourth moment.
        \item[4.] Gamma regression: $Y>0$ almost surely, $X$ is Gaussian or sub-Gaussian, $Y$ has a finite fourth moment, and $\mathbb E\{(\log Y)^2\}<\infty$.
    \end{enumerate}
\end{proposition}
\begin{remark}
\begin{itemize}
    \item[i)] A random vector $X$ has a finite $p$th moment when $\mathbb E(\lVert X\rVert^p)<\infty$, where $\lVert\cdot\rVert$ is the Euclidean norm.
    \item[ii)] The same conditions hold when both the covariate vector and the response have compact support compatible with the relevant response space. For Gamma regression, the compact response support must be bounded away from zero so that $\log Y$ is bounded.
    \item[iii)] Throughout, a random vector $X$ is called sub-Gaussian if there is a constant $K<\infty$ such that
    \[
    \mathbb E\exp\{t u^\top(X-\mathbb EX)\}
    \le \exp\left(\frac{K^2t^2\|u\|_2^2}{2}\right)
    \quad\text{for all }u\in\mathbb R^d\text{ and }t\in\mathbb R.
    \]
    In fixed dimension this implies $\mathbb E\exp\{s\|X\|_1\}<\infty$ for every $s>0$, as used in the Poisson and Gamma envelope arguments.
\end{itemize}
\end{remark}

\begin{bluechange}
Let $\mathcal B_k$ denote the compact slope constraint set in candidate $k$, let
\[
B_{\mathbb{U}}=\sup_{u\in\mathbb{U}}|u|,
\qquad B_k=\sup_{\beta\in\mathcal B_k}\|\beta\|_2,
\qquad \widetilde x=(1,x^\top)^\top,
\]
and define $A_k(x)=B_{\mathbb{U}}+B_k\|x\|_2$.

\begin{proposition}\label{Prop: intercept envelopes}
For $\theta=(\beta_0,\beta)$ and the Euclidean parameter norm, the four losses admit the following pointwise Lipschitz moduli on $\mathbb{U}\times\mathcal B_k$:
\begin{align*}
\mathcal G_k^{\mathrm{lin}}(x,y)&=2\{|y|+A_k(x)\}\|\widetilde x\|_2,\\
\mathcal G_k^{\mathrm{logit}}(x,y)&=\|\widetilde x\|_2,\\
\mathcal G_k^{\mathrm{Pois}}(x,y)&=\{\exp(A_k(x))+y\}\|\widetilde x\|_2,\\
\mathcal G_k^{\mathrm{Gamma}}(x,y)&=\nu\{1+y\exp(A_k(x))\}\|\widetilde x\|_2.
\end{align*}
Under the corresponding conditions of Proposition \ref{Prop: Sufficient conditions}, the square of the relevant envelope is integrable. Thus A3 holds with the intercept included; A2 also holds by taking any fixed $u_0\in\mathbb{U}$ and $\beta=0$.
\end{proposition}

For joint estimation of the Gamma shape, let $\mathbb{V}=[\nu_-,\nu_+]$ be a fixed compact interval with $0<\nu_-<\nu_+<\infty$, and use the full negative log-likelihood, up to model-independent constants,
\begin{equation}\label{Eqs: Gamma estimated loss}
\ell(x,y;\beta_0,\beta,\nu)
=\log\Gamma(\nu)-\nu\log\nu+\nu\eta-(\nu-1)\log y+\nu y\exp(-\eta),
\qquad \eta=\beta_0+x^\top\beta.
\end{equation}
\begin{proposition}\label{Prop: estimated Gamma shape}
Assume $Y>0$ almost surely, $X$ is Gaussian or sub-Gaussian, $\mathbb E(Y^4)<\infty$, and $\mathbb E\{(\log Y)^2\}<\infty$. If the candidate parameter spaces are $\mathbb{U}\times\mathcal B_k\times\mathbb{V}$, with all factors compact and the slope sets nested, then the loss in (\ref{Eqs: Gamma estimated loss}) satisfies A1--A3. Consequently, every PanIC penalty satisfying B1 and B2 consistently selects the smallest optimal constraint level while $\nu$ is estimated jointly within each candidate model.
\end{proposition}
\end{bluechange}

\begin{bluechange}
Let $D_{k,n}$ be any nonnegative measurable empirical complexity score, and define its cumulative maximum by
\begin{equation}\label{Eqs: cumulative complexity}
\widetilde D_{1,n}=D_{1,n},\qquad
\widetilde D_{k,n}=\max\{\widetilde D_{k-1,n},D_{k,n}\}
=\max_{1\le j\le k}D_{j,n},\quad k\ge2.
\end{equation}
For a fixed grid $C_1<\cdots<C_m$, a nonnegative strictly increasing function $\psi$ on the grid, and constants $c_0>0$, $\kappa_D\ge0$, and $\varepsilon>0$, define, for $n\ge2$,
\begin{equation}\label{Eqs: Modified BIC}
\widehat K_n=\min\left\{\argmin_{k\in[m]}\left[
R_n(\widehat\beta_{0,k,n},\widehat\beta_{k,n})
+\frac{\log n}{n}\{c_0+\kappa_D\widetilde D_{k,n}+\varepsilon\psi(C_k)\}
\right]\right\}.
\end{equation}

\begin{remark}
The criterion in (\ref{Eqs: Modified BIC}) is BIC-like on the mean-squared-error scale. To make the scaling explicit, write
\[
\mathrm{RSS}_k=\sum_{i=1}^n\left(Y_i-\widehat\beta_{0,k,n}-X_i^\top\widehat\beta_{k,n}\right)^2,
\qquad R_{k,n}=\frac{\mathrm{RSS}_k}{n}.
\]
For Gaussian errors of known variance $\sigma^2$, let $d_k$ denote the candidate-specific model dimension used in the BIC. The likelihood-based BIC of \citet{schwarz_estimating_1978} is, up to a model-independent constant, $\mathrm{RSS}_k/\sigma^2+d_k\log n$. Hence minimising it is equivalent to minimising
\begin{equation}\label{Eqs: known variance BIC}
R_{k,n}+\sigma^2d_k\frac{\log n}{n}.
\end{equation}
Thus a unit coefficient on the MSE-scale complexity term is exact only when $\sigma^2=1$ or the squared loss has been divided by $\sigma^2$. If $R_{k,n}>0$ and the variance is estimated separately in each model, profiling it out yields, up to a model-independent constant, $n\log(\mathrm{RSS}_k/n)+d_k\log n$, or equivalently
\begin{equation}\label{Eqs: profiled BIC}
\log(R_{k,n})+d_k\frac{\log n}{n}.
\end{equation}
\end{remark}

\begin{proposition}\label{Prop: generic BIC penalty}
Suppose $m<\infty$ is fixed, $(r_n)$ is a deterministic positive sequence, and
\[
\max_{1\le k\le m}D_{k,n}=O_{\mathbb P}(r_n),
\qquad \frac{r_n\log n}{n}\longrightarrow0.
\]
For $n\ge2$, the penalty in (\ref{Eqs: Modified BIC}) is strictly positive, is $o_{\mathbb P}(1)$ uniformly over $k$, and satisfies $\mathrm{B2}^*$. Define the penalty at $n=1$ to be any positive value. Consequently, whenever A1--A3 and C1--C5 hold, the selector in (\ref{Eqs: Modified BIC}) is consistent.
\end{proposition}

The consistency property only requires a nonnegative bounded or slowly growing score $D_{k,n}$; a more specific degrees-of-freedom interpretation is model-specific. For centred Lasso least squares, the number of non-zero slope parameters is an unbiased estimate of the degrees of freedom of the slope fit at a fixed penalty multiplier under the Gaussian-model and full-rank design conditions of \citet{zou_degrees_2007}. An unrestricted unpenalised intercept contributes one additional degree of freedom, common to all candidates. For the compact-intercept fits considered here, or at a data-dependent multiplier or a prescribed constraint radius, the count is used only as a generic complexity score. For Ridge least squares with an unpenalised intercept, let $\check X$ denote the column-centred design matrix and let $\check y$ be the centred response vector, with $(\check y)_i=y_i-n^{-1}\sum_{j=1}^n y_j$. For a fixed $\lambda>0$, the penalised objective $n^{-1}\|\check y-\check X\beta\|_2^2+\lambda\|\beta\|_2^2$ has centred slope smoother trace
\begin{equation}\label{Eqs: Ridge df}
D^{\mathrm{ridge}}(\lambda)
=\operatorname{tr}\{\check X(\check X^\top \check X+n\lambda I)^{-1}\check X^\top\}.
\end{equation}
At $\lambda=0$, the inverse is replaced by the Moore--Penrose pseudoinverse, giving $\operatorname{rank}(\check X)$. Evaluating this trace at a data-dependent multiplier supplies a generic score, not automatically the degrees of freedom of the complete selection procedure.
For Elastic-net, a non-zero parameter count may be used as a bounded generic complexity score, but it is not described here as an unbiased degrees-of-freedom estimator.

\begin{remark}
For any common random or deterministic scalar $a_n$ and any $c_0$ not depending on $k$,
\[
\argmin_k\{v_{k,n}+a_n(c_0+q_k)\}
=\argmin_k\{v_{k,n}+a_nq_k\}.
\]
Thus $c_0$ cancels from every pairwise criterion difference and cannot change the selected model; it only enforces strict positivity when the smallest threshold is zero.
\end{remark}

Because coefficient-norm thresholds depend on predictor units, all continuous predictors should be centred and scaled, with the intercept left unpenalised. For Gaussian regression, the response or loss should likewise be put on a documented dimensionless scale. A dimensionless threshold transform on a deterministic interval $[a,b]$ is $\psi(C)=(C-a)/(b-a)$; on a fixed grid with $m\geq2$, one may instead use the rank $(k-1)/(m-1)$. Any such strictly increasing transform preserves the consistency argument.
\end{bluechange}

\begin{proposition}
    \label{Prop: BIC-like consistency}
    In the context of linear regression, let $(X, Y)$ be a random element on the measurable space $(\mathbb{R}^d \times \mathbb{R}, \mathcal{B}(\mathbb{R}^d)\otimes \mathcal{B}(\mathbb{R}))$ with law $\Pi$. Assume that both $X$ and $Y$ have finite fourth moments, and that the covariance matrix $\Sigma$ of $X$ is positive definite. That is,
    \[\Sigma = \mathbb{E}\{(X - \mathbb{E}\{X\})(X - \mathbb{E}\{X\})^\top\} \qquad \text{satisfies} \qquad u^\top \Sigma u > 0\]
    for every $u \in \mathbb{R}^d \backslash \{0\}$. Then conditions C1--C5 hold for the regularised linear least-squares problem, so every criterion satisfying Proposition \ref{Prop: generic BIC penalty}, including the Lasso non-zero-parameter-count version, is consistent. When an intercept is included, write $\tilde{X}=(1,X^\top)^\top$ and
    \[
    \tilde{\Sigma}=\mathbb{E}(\tilde{X}\tilde{X}^\top)=\begin{bmatrix}1 & \mu^\top\\ \mu & \Sigma+\mu\mu^\top\end{bmatrix}.
    \]
    Here $\mu=\mathbb{E}(X)$; the assumption that $\Sigma$ is positive definite implies that $\tilde{\Sigma}$ is positive definite, which is the matrix appearing in the Hessian with respect to $(\beta_0,\beta)$. The compact interval constraining the intercept is common to all candidates and is polyhedral; taking its Cartesian product with each slope constraint preserves the second-order regularity verified in the proof.
\end{proposition}

To allow hypotheses indexed by a compact interval, we use a continuous, compact-valued correspondence. A correspondence $\varphi$ from a set $\mathbb{X}$ to a set $\mathbb{Y}$ assigns to each $x\in\mathbb{X}$ a subset of $\mathbb{Y}$ \citep[Definition 17.1]{aliprantis_infinite_2006}.

\begin{theorem}\label{Thm: Infinite selection}
Let $(X_{i})_{i\in\mathbb{N}}$ be an i.i.d. sequence taking values in a set $\mathbb{X}\subset\mathbb{R}^{d}$, and let $(\mathcal{H}_{k})_{k\in\left[a,b\right]}$ be hypothesis spaces of the form
\[
\mathcal{H}_{k}=\left\{ h(\cdot;\beta):\mathbb{X}\to\mathbb{R}, \beta\in\mathbb{T}_{k}\subset\mathbb{R}^{q}\right\},
\]
for some positive integer $q$, where $a<b$. Let $\mathcal{L}:k\mapsto\mathbb{T}_{k}$ be a continuous correspondence with \blue{nonempty} compact values, and assume that the parameter spaces are increasing in the index, so that $\mathbb T_k\subseteq\mathbb T_{k'}$ whenever $k<k'$. Let $\ell:\mathbb{X}\times\mathbb{R}^{q}\to\mathbb{R}$ be a loss function and define
\[
R_{n}(\beta)=\frac{1}{n}\sum_{i=1}^{n}\ell(X_{i};\beta),\qquad r(\beta)=\mathbb{E}\{\ell(X;\beta)\}.
\]
Assume that, for each $k\in\left[a,b\right]$, the restriction of $\ell$ to $\mathbb{T}_{k}$ satisfies A1--A3. Define
\[
\mathcal{K}=\argmin_{k\in\left[a,b\right]}\left\{ \min_{\beta\in\mathbb{T}_{k}}r(\beta)\right\},\qquad k^{*}=\min_{k\in\mathcal{K}}k.
\]
Let $(P_{k,n})_{k\in\left[a,b\right]}$ be non-negative measurable penalties. Assume that, for each $n\in\mathbb{N}$, the sample path $k\mapsto P_{k,n}(\omega)$ is continuous and strictly increasing on $\left[a,b\right]$ for almost every $\omega$, and that B1 holds. Define the associated selector
\[
\hat{K}_{n}=\min\left\{\argmin_{k\in\left[a,b\right]}\left\{ \min_{\beta\in\mathbb{T}_{k}}R_{n}(\beta)+P_{k,n}\right\}\right\}.
\]
If $P_{k,n}$ satisfies either B2, or $\mathrm{B2}^{*}$ and C1--C5 (with $k$ ranging over $\left[a,b\right]$), then $\hat{K}_{n}$ is consistent. That is,
\begin{equation}
\plim_{n\to\infty}\hat{K}_{n}=k^{*}.\label{Eqs: convergence in probability}
\end{equation}
\end{theorem}

For a continuum interval $[a,b]$ with $0\le a<b$, define the regression parameter space $\mathbb{T}_k$ of the correspondence $\mathcal{L}:k\mapsto\mathbb{T}_k$ by
\begin{equation}\label{Eqs: Parameter Space}
    \mathbb{T}_k = \begin{cases}
        \{(\beta_0, \beta) \, |\, \beta_0 \in \mathbb{U}, \beta\in \mathbb{R}^d: \lVert\beta\rVert_1 \leq k\}, & \text{for Lasso regression}\\
        \{(\beta_0, \beta) \, |\, \beta_0 \in \mathbb{U}, \beta\in \mathbb{R}^d: \lVert\beta\rVert_2^2 \leq k\}, & \text{for Ridge regression}\\
        \{(\beta_0, \beta) \, |\, \beta_0 \in \mathbb{U}, \beta\in \mathbb{R}^d: \alpha\lVert\beta\rVert_1 + (1-\alpha)\lVert\beta\rVert_2^2 \leq k\}, & \text{for Elastic-net regression}.
    \end{cases}
\end{equation}
\begin{bluechange}
\begin{lemma}\label{Lemma: constraint correspondence continuity}
For each of the three definitions in (\ref{Eqs: Parameter Space}), the correspondence $k\mapsto\mathbb T_k$ is nonempty, compact-valued, increasing, and continuous on $[a,b]$.
\end{lemma}
\begin{proof}
Write $\mathbb T_k=\mathbb U\times\mathcal B_k$, where $\mathcal B_k=\{\beta:g(\beta)\le k\}$ and $g$ is the Lasso, Ridge, or Elastic-net constraint functional. In all three cases, $g$ is continuous, coercive, and satisfies $g(0)=0$. Hence every $\mathcal B_k$ is nonempty and compact, and the sets are increasing in $k$.

Let $k_n\to k$. If $\beta_n\in\mathcal B_{k_n}$ and $\beta_n\to\beta$, continuity of $g$ gives $g(\beta)\le k$, which proves upper hemicontinuity; coercivity supplies the local uniform boundedness needed for the sequential argument. For lower hemicontinuity, fix $\beta\in\mathcal B_k$. If $k=0$, then $\beta=0$ and one may take $\beta_n=0$. If $k>0$, set $t_n=\min\{1,k_n/k\}$ and $\beta_n=t_n\beta$. For $0\le t\le1$, each of the three functionals satisfies $g(t\beta)\le t g(\beta)$, so $g(\beta_n)\le k_n$ and $\beta_n\to\beta$. Thus $k\mapsto\mathcal B_k$ is continuous. Taking its Cartesian product with the fixed compact interval $\mathbb U$ preserves continuity and compact-valuedness.
\end{proof}
\end{bluechange}

Theorem \ref{Thm: Infinite selection} and Lemma \ref{Lemma: constraint correspondence continuity} give the following result.

\begin{corollary}\label{Cor: continuum PanIC GLM}
Under the assumptions of Theorem \ref{Thm: Infinite selection}, let $(\mathcal H_k)_{k\in[a,b]}$ be the family of hypothesis spaces for a linear, logistic, Poisson, or Gamma regression problem described in Table \ref{Table: Regression}, with parameter spaces as in (\ref{Eqs: Parameter Space}). Suppose the corresponding conditions in Proposition \ref{Prop: Sufficient conditions} hold and $P_{k,n}$ satisfies B1 and B2. Then the selector in Theorem \ref{Thm: Infinite selection} is consistent in the sense of (\ref{Eqs: convergence in probability}).
\end{corollary}

\begin{corollary}\label{Cor: continuum BIC linear}
In the linear-regression setting, suppose the conditions of Proposition \ref{Prop: BIC-like consistency} hold. Then the conclusion of Theorem \ref{Thm: Infinite selection} applies to any continuous strictly increasing BIC-like penalty on $[a,b]$ that satisfies B1 and $\mathrm{B2}^*$. In particular, for any fixed continuous strictly increasing $q:[a,b]\to\mathbb R_{>0}$, the penalty $P_{k,n}=q(k)\log(n)/n$ for $n\ge2$, together with a positive continuous strictly increasing convention at $n=1$, such as $P_{k,1}=q(k)$, satisfies these conditions.
\end{corollary}

\begin{bluechange}
The deterministic continuum theorem does not by itself cover a random list of thresholds. The following transfer result supplies a sufficient condition. Write
\[
v(c)=\min_{\theta\in\mathbb T_c}r(\theta),\qquad
v_n(c)=\min_{\theta\in\mathbb T_c}R_n(\theta),
\]
and let $c^*$ be the smallest minimiser of $v$.

\begin{theorem}\label{Thm: random finite grid}
Let $(\mathbb T_c)_{c\in[a,b]}$ be increasing, so that $\mathbb T_c\subseteq\mathbb T_{c'}$ whenever $c<c'$, and let $\mathcal G_n\subset[a,b]$ be a nonempty finite random set, possibly constructed from the same data used in $R_n$. Suppose there are measurable $c_n^+\in\mathcal G_n$ and nonnegative $\delta_n=o_{\mathbb P}(1)$ such that
\begin{equation}\label{Eqs: one sided coverage}
c^*\le c_n^+\le c^*+\delta_n.
\end{equation}
Let $P_n(c)$ be nonnegative and increasing in $c$, with $P_n(b)=o_{\mathbb P}(1)$. Assume that, for some deterministic $a_n\to\infty$,
\begin{equation}\label{Eqs: random grid empirical gap}
a_n\{v_n(b)-v_n(c^*)\}=O_{\mathbb P}(1),
\end{equation}
and, for every $\epsilon>0$ such that $c^*+\epsilon\le b$,
\begin{equation}\label{Eqs: random grid penalty separation}
a_n\{P_n(c^*+\epsilon)-P_n(c^*+\epsilon/2)\}
\xlongrightarrow{\mathbb P}\infty.
\end{equation}
Finally, for every $\epsilon>0$ with $c^*-\epsilon\ge a$, assume that
\begin{equation}\label{Eqs: random grid pointwise convergence}
v_n(c^*-\epsilon)\xlongrightarrow{\mathbb P}v(c^*-\epsilon),
\qquad
v_n(c^*)\xlongrightarrow{\mathbb P}v(c^*).
\end{equation}
Assume that the minimum tie-broken minimiser below is measurable, and define
\[
\widehat c_n^{\mathcal G}
=\min\left\{\argmin_{c\in\mathcal G_n}\{v_n(c)+P_n(c)\}\right\}.
\]
Then $\widehat c_n^{\mathcal G}\to_{\mathbb P}c^*$.
\end{theorem}

For PanIC one may take $a_n=\sqrt n$, while for a BIC-like criterion under C1--C5 one may take $a_n=n$. To express the one-sided coverage condition directly, define
\[
\Delta_n^+=\inf\{c-c^*:c\in\mathcal G_n,\ c\ge c^*\},
\]
with the infimum of the empty set defined as $+\infty$. Condition (\ref{Eqs: one sided coverage}) implies $\Delta_n^+=o_{\mathbb P}(1)$ and hence
\[
\mathbb P\{\mathcal G_n\cap[c^*,c^*+\eta]\ne\varnothing\}\longrightarrow1
\quad\text{for every }\eta>0.
\]
Conversely, the displayed probability condition is equivalent to $\Delta_n^+=o_{\mathbb P}(1)$, but it permits an empty upper grid on exceptional events. The theorem uses the stronger explicit formulation in (\ref{Eqs: one sided coverage}), which supplies an actual measurable candidate on every sample; a theorem formulated only in probability would also need a convention on exceptional samples. Thus arbitrary regularisation-path knots are not automatically covered. The condition may be verified for a particular path construction or enforced by augmentation with a deterministic mesh. For example, let $(J_n)$ be a sequence of positive integers with $J_n\to\infty$, set $h_n=(b-a)/J_n$, and define
\[
\mathcal D_n=\{a+jh_n: j=0,\ldots,J_n\}.
\]
For any fixed $c^*\in[a,b]$, the point $c_n^+=\min\{d\in\mathcal D_n:d\ge c^*\}$ exists and satisfies $0\le c_n^+-c^*\le h_n\to0$. Hence, if $\mathcal P_n\subset[a,b]$ denotes any finite data-derived candidate set, the augmented set $\mathcal G_n=\mathcal P_n\cup\mathcal D_n$ satisfies the one-sided coverage condition deterministically. The remaining assumptions of Theorem \ref{Thm: random finite grid}, including the empirical-risk and penalty-separation conditions, must still be verified.

When the candidate grid is fixed, every adjacent transformed gap $\psi(C_{k+1})-\psi(C_k)$ is a positive constant, so the usual pairwise checks of B2 and $\mathrm{B2}^*$ apply. If the grid itself changes with $n$ and the target is exact recovery of a particular grid index, adjacent penalties must not become too close relative to the stochastic scale. For an ordered grid $C_{1,n}<\cdots<C_{m_n,n}$, define
\[
\Delta_{\psi,n}=\min_{1\le k<m_n}\{\psi(C_{k+1,n})-\psi(C_{k,n})\}.
\]
For a PanIC shape proportional to $\psi(C)\sqrt{\log n/n}$, the sufficient adjacent-gap condition is $\sqrt{\log n}\,\Delta_{\psi,n}\to\infty$; for a BIC-like shape proportional to $\psi(C)\log n/n$, it is $\log n\,\Delta_{\psi,n}\to\infty$. When $\psi(C)=C$, these are conditions on the raw threshold gaps. They ensure only the required penalty separation; exact index consistency also requires the remaining model-selection assumptions and suitable uniform empirical-risk control for the changing collection. By contrast, threshold consistency under Theorem \ref{Thm: random finite grid} does not require a global minimum adjacent gap. It compares candidates separated from $c^*$ by each fixed positive distance and separately requires a candidate approaching $c^*$ from above.

\begin{remark}\label{Rem: multiplier continuum random grid}
The compact-range multiplier argument is compatible with both continuum and random-grid selection. If a base penalty $Q_n(c)$ satisfies the relevant vanishing and fixed-separation conditions and $\widehat\kappa_n>0$ almost surely and lies with probability tending to one in $[\underline\kappa,\overline\kappa]$, where $0<\underline\kappa\le\overline\kappa<\infty$, then $P_n(c)=\widehat\kappa_nQ_n(c)$ preserves those conditions by the same probability bounds as Proposition \ref{Prop: uniform kappa}. Positive multiplication also preserves continuity and any strict monotonicity of the penalty sample paths. For Theorem \ref{Thm: random finite grid}, it preserves $P_n(b)=o_{\mathbb P}(1)$ and the fixed-distance separation condition. The calibration does not, however, establish the one-sided coverage condition or the empirical-risk assumptions; those properties concern the candidate set and the fitted risks rather than the multiplier.
\end{remark}
\end{bluechange}

\subsection{Computation}
\label{Sect: Computations}

PanIC and, for regularised linear regression, the BIC-like criterion in (\ref{Eqs: Modified BIC}) can be computed on a finite ordered candidate set satisfying the preceding finite-grid conditions.

Suppose $f:\mathbb R^d\to\mathbb R$ is convex and $g:\mathbb R^d\to\mathbb R_{\geq0}$ is nonnegative and convex. Consider the following constrained optimisation problem:
\begin{align}
\begin{split}\min\quad & \,f(\beta)\\
\mathrm{s.t.}\quad & \beta\in\mathbb{R}^{d}:g(\beta)\leq C,
\end{split}
\label{Eqs: General constrained problem}
\end{align}
for some value $C\in\mathbb{R}_{+}$. Assume that a solution exists and that the convex problem satisfies a constraint qualification such as Slater's condition. The Karush--Kuhn--Tucker conditions then imply that each constrained optimum admits a multiplier $\lambda\ge0$ and minimises the penalised objective
\begin{align}
\begin{split}\min\quad & \,f(\beta)+\lambda g(\beta)\\
\mathrm{s.t.}\quad & \beta\in\mathbb{R}^{d}.
\end{split}
\label{Eqs: General unconstrained problem}
\end{align}
Conversely, a minimiser $\beta_\lambda$ of the penalised problem is a constrained minimiser at the level $C_\lambda=g(\beta_\lambda)$. These implications do not yield a one-to-one correspondence between $C$ and $\lambda$: multipliers may be nonunique and the fitted constraint value may have plateaus \citep{kloft_lp-norm_2011}. Under the additional uniqueness assumptions below, define the fitted-radius map $\rho(\lambda)=g(\beta_\lambda)$, where $\beta_\lambda$ is the unique penalised minimiser. Lemmas \ref{Lemma: continuity} and \ref{Lemma: monotonicity} establish the continuity and monotonicity needed for root finding. Although these lemmas are written without an intercept, their proofs apply on $\mathbb U\times\mathbb R^d$ after replacing each sublevel set by the compact set $\mathbb U\times\{\beta:g(\beta)\le c\}$ and penalising only $g(\beta)$.
\begin{lemma}\label{Lemma: continuity}
Suppose $f:\mathbb{R}^{d}\to\mathbb{R}$ is a strictly convex
function and $g:\mathbb{R}^{d}\to\mathbb{R}_{\geq0}$ is a non-negative
convex function with non-empty and bounded sub-level sets $L_{c}=\{\beta\in\mathbb{R}^{d}\,|\,g(\beta)\leq c\}$,
for each $c\in\mathbb{R}_{\geq0}$. Assume that, for each $\lambda\in\mathbb{R}_{\geq0}$, the function
$\beta\mapsto f\left(\beta\right)+\lambda g\left(\beta\right)$ admits a unique minimiser on $\mathbb{R}^{d}$. Then, the fitted-radius map $\rho:\mathbb{R}_{\geq0}\to\mathbb{R}_{\geq0}$,
given by
\[
\rho(\lambda)=g(\beta_{\lambda}),
\]
is a continuous function, where
\[
\beta_{\lambda}=\underset{\beta\in\mathbb{R}^{d}}{\arg\min}\ f\left(\beta\right)+\lambda g\left(\beta\right).
\]
\end{lemma}

\begin{lemma}\label{Lemma: monotonicity} Let $\lambda_{1}$ and
$\lambda_{2}$ be non-negative real numbers, with $\lambda_{2}>\lambda_{1}$.
Under the same assumptions as those of Lemma \ref{Lemma: continuity},
it holds that $g(\beta_{\lambda_{2}})\leq g(\beta_{\lambda_{1}})$.
Further, if $g(\beta_{\lambda_{2}})<g(\beta_{\lambda_{1}})$, then
$f(\beta_{\lambda_{2}})>f(\beta_{\lambda_{1}})$. Else, if $g(\beta_{\lambda_{2}})=g(\beta_{\lambda_{1}})$,
then $\beta_{\lambda_{2}}=\beta_{\lambda_{1}}$. \end{lemma}

For a candidate radius $C$, the constrained regression problem is
\begin{align}
\begin{split}\min\quad & \,R_{n}(\beta_{0},\beta)=\frac{1}{n}\sum_{i=1}^{n}\ell(x_{i},y_{i};\beta_{0},\beta)\\
\mathrm{s.t.}\quad & \beta_{0}\in\mathbb{U},\,\beta\in\mathbb{R}^{d}:g(\beta)\leq C,
\end{split}
\label{Eqs: Constrained problem}
\end{align}
for some $C\in\mathbb R_{\geq0}$. Here, $g$ is the $l_1$ norm, the squared $l_2$ norm, or their Elastic-net combination. For the convex computational formulation, $\mathbb U$ is taken to be a compact interval and the relevant intercept solution is assumed to lie in its interior. For $C>0$, the zero slope vector together with an interior intercept is strictly feasible, so Slater's condition holds. If the penalised minimiser is unique for every $\lambda\ge0$, Lemmas \ref{Lemma: continuity} and \ref{Lemma: monotonicity} permit bisection or another bracketed root search for an active constraint. \blue{If uniqueness fails, the constrained problem may instead be solved directly. A deterministic choice among penalised minimisers does not by itself give the continuity required for bisection.}

\begin{bluechange}
For an active constraint, the continuity and monotonicity in Lemmas \ref{Lemma: continuity} and \ref{Lemma: monotonicity} permit a bracketed search for $g(\beta_\lambda)=C$. For the slope penalties considered here, $g(0)=0$ and
\[
0\le g(\beta_\lambda)
\le \frac{f(0)-\min_\beta f(\beta)}{\lambda},\qquad \lambda>0.
\]
Thus the fitted radius tends to zero as $\lambda\to\infty$, so a bracket exists for every $0<C<g(\beta_{\lambda=0})$ under the lemma's assumptions. With an intercept, the same argument uses a fixed feasible intercept at zero slope as the reference point. An inactive constraint returns the unpenalised fit, and $C=0$ returns the zero-slope fit.
\end{bluechange}

\begin{bluechange}
The numerical study maps an unpenalised-intercept \texttt{glmnet} path to the fixed radius grid; this approximates the exact constrained minimisers used in the theory. Agreement with the compact-intercept theory requires the relevant population intercepts to lie in the interior of a fixed compact interval and the intercept estimates to be consistent over the finite candidate collection. The code repository cited in Appendix \ref{App: reproducibility} gives the implementation and failure-handling details.

The empirical comparisons cover Gaussian, logistic, and Poisson regression. Gamma regression is treated analytically in Proposition \ref{Prop: estimated Gamma shape} but is not included in the comparative simulations.
\end{bluechange}

Algorithm \ref{Algo: PanIC Lasso} computes the constrained fits once and then selects the smallest minimiser of $\widehat v_{k,n}+P_{k,n}$. The penalty input may be the PanIC penalty or, for linear regression, the BIC-like penalty in (\ref{Eqs: Modified BIC}).

\begin{algorithm}[H]
\small
\caption{Computing the Lasso solution with PanIC.}
\label{Algo: PanIC Lasso} \begin{algorithmic}[1] \Require a data sample $(x_{i},y_{i})_{i\in\left[n\right]}$, a strictly increasing sequence
$(C_{k})_{k\in\left[m\right]}$, a function $h(\cdot;\beta_{0},\beta):\mathcal{X}\to\mathcal{M}$,
a loss function $\ell(x,y;\beta_{0},\beta)$,
a penalty term $P_{k,n}$ (for example, the PanIC penalty, or the
modified BIC penalty term in (\ref{Eqs: Modified BIC})). \State
Formulate a sequence of hypotheses $(\mathcal{H}_{k})_{k\in\left[m\right]}$
with
\[
\mathbb{T}_{k}\gets\{(\beta_{0},\beta)\,|\,\beta_{0}\in\mathbb{U},\,\beta\in\mathbb{R}^{d}:\lVert\beta\rVert_{1}\leq C_{k}\} \quad \text{and} \quad \mathcal{H}_{k}\gets\{h(\cdot;\beta_{0},\beta)\,|\,(\beta_{0},\beta)\in\mathbb{T}_{k}\}.
\]
Associated with each hypothesis, define the optimisation problem
\begin{align}
\begin{split}\min\quad & \,R_{n}(\beta_{0},\beta)\\
\mathrm{s.t.}\quad & (\beta_{0},\beta)\in\mathbb{T}_{k}.
\end{split}
\label{Eqs: Constrained problem algo}
\end{align}
\For{$k\in\left[m\right]$} \State Solve (\ref{Eqs: Constrained problem algo})
directly, or via the corresponding Lagrangian problem when the root-search conditions above hold; return the unpenalised fit when the constraint is inactive.
\State Denote the optimal value and optimal solution of (\ref{Eqs: Constrained problem algo}) by $\widehat v_{k,n}$ and $(\hat{\beta}_{0,k,n},\hat{\beta}_{k,n})$, respectively.
\EndFor \State Compute the set of minimising indices $\hat{\mathcal{K}}$, and then compute the PanIC estimate, where
\[
\hat{\mathcal{K}}\gets\argmin_{k\in\left[m\right]}\left\{ \widehat v_{k,n}+P_{k,n}\right\} \qquad \text{and} \qquad \hat{K}_{n}\gets\min_{k\in\hat{\mathcal{K}}}\,k.
\]
\Return $\hat{K}_{n}$, $(\hat{\beta}_{0,\hat{K}_{n},n},\hat{\beta}_{\hat{K}_{n},n})$
\end{algorithmic}
\end{algorithm}

\subsection{Second-order asymptotics for constrained linear regression}

\label{Sect: Asymptotic-distribution}

We derive the second-order law for constrained least squares. A binding constraint generally replaces the usual OLS Gaussian limit by the solution of a random quadratic programme over the critical cone.

We work under the following setup. Let $(X_{i},Y_{i})_{i\in\mathbb{N}}$
be i.i.d. with $X_{i}\in\mathbb{R}^{d}$. We estimate an intercept
$\beta_{0}\in\mathbb{R}$ and slopes $\beta\in\mathbb{R}^{d}$, and
define the squared loss
\[
\ell(x,y;\beta_{0},\beta)=\left(y-\beta_{0}-x^{\top}\beta\right)^{2},\qquad R_{n}(\beta_{0},\beta)=\frac{1}{n}\sum_{i=1}^{n}\ell(X_{i},Y_{i};\beta_{0},\beta),\qquad r(\beta_{0},\beta)=\mathrm{E}\left\{ \ell(X,Y;\beta_{0},\beta)\right\} .
\]

For each $k$, let $\mathbb{T}_{k}\subset\mathbb{R}^{d+1}$ be the
parameter space defined in (\ref{Eqs: Parameter Space}), with constraint
level $C_{k}$ in place of $k$. That is, for a fixed compact interval
$\mathbb{U}\subset\mathbb{R}$ as in (\ref{Eqs: Lasso Ridge}),
\begin{align*}
\textnormal{(Lasso)}\quad & \mathbb{T}_{k}=\left\{ (\beta_{0},\beta)\in\mathbb{R}\times\mathbb{R}^{d}:\beta_{0}\in\mathbb{U},\ \lVert\beta\rVert_{1}\le C_{k}\right\} ,\\
\textnormal{(Ridge)}\quad & \mathbb{T}_{k}=\left\{ (\beta_{0},\beta)\in\mathbb{R}\times\mathbb{R}^{d}:\beta_{0}\in\mathbb{U},\ \lVert\beta\rVert_{2}^{2}\le C_{k}\right\} ,\\
\textnormal{(Elastic-net)}\quad & \mathbb{T}_{k}=\left\{ (\beta_{0},\beta)\in\mathbb{R}\times\mathbb{R}^{d}:\beta_{0}\in\mathbb{U},\ \alpha\lVert\beta\rVert_{1}+\left(1-\alpha\right)\lVert\beta\rVert_{2}^{2}\le C_{k}\right\} ,
\end{align*}
with $(C_{k})$ strictly increasing.

Define the population and empirical constrained minimisers
\[
(\beta_{0,k}^{*},\beta_{k}^{*})=\argmin_{(\beta_{0},\beta)\in\mathbb{T}_{k}}r(\beta_{0},\beta),\qquad(\hat{\beta}_{0,k,n},\hat{\beta}_{k,n})\in\argmin_{(\beta_{0},\beta)\in\mathbb{T}_{k}}R_{n}(\beta_{0},\beta),
\]
and the corresponding optimal values
\[
v_{k}=r(\beta_{0,k}^{*},\beta_{k}^{*}),\qquad\hat{v}_{k,n}=R_{n}(\hat{\beta}_{0,k,n},\hat{\beta}_{k,n}).
\]
\begin{bluechange}
Throughout this subsection, fix $k$ and write $\theta_k^*=(\beta_{0,k}^*,\beta_k^{*\top})^\top$. We assume that $\theta_k^*$ is the unique population minimiser, that $\mathbb T_k$ is second-order regular at $\theta_k^*$, and that $\beta_{0,k}^*$ lies in the interior of $\mathbb U$. We also assume
\[
\mathbb E\left(\|\widetilde X\|_2^4\right)<\infty,
\qquad
\mathbb E(Y^4)<\infty,
\qquad
H=2\mathbb E(\widetilde X\widetilde X^\top)\ \text{is positive definite},
\]
where $\widetilde X=(1,X^\top)^\top$. The fourth-moment assumptions imply $\mathbb E\{Y^2\|\widetilde X\|_2^2\}<\infty$ by the Cauchy--Schwarz inequality and hence give a finite score covariance. Moreover, the squared loss is a quadratic polynomial in $\theta$: its empirical coefficients are averages of $Y^2$, $Y\widetilde X$, and $\widetilde X\widetilde X^\top$. Their joint multivariate central limit theorem therefore yields $\sqrt n(R_n-r)$ convergence in the function-and-gradient supremum norm on the compact set $\mathbb T_k$, while the law of large numbers gives convergence of the empirical Hessian. Together with positive definiteness of $H$, the exact quadratic expansion, and second-order regularity, these facts verify the stochastic-optimisation conditions used in \citet[Theorem 4.4]{shapiro_statistical_2000}.
\end{bluechange}

Let $\mathcal{T}_{\mathbb{T}_{k}}(\beta_{0,k}^{*},\beta_{k}^{*})$ denote the
tangent cone of $\mathbb{T}_{k}$ at $(\beta_{0,k}^{*},\beta_{k}^{*})$,
and let $\mathcal{T}_{\mathbb{T}_{k}}^{2}\left((\beta_{0,k}^{*},\beta_{k}^{*}),d\right)$
denote the second-order tangent set in direction $d$, as in \citet[Section 4]{shapiro_statistical_2000}.
Define the critical cone
\begin{equation}
\mathcal{C}_{k}(\beta_{0,k}^{*},\beta_{k}^{*})=\left\{ d\in \mathcal{T}_{\mathbb{T}_{k}}(\beta_{0,k}^{*},\beta_{k}^{*}):d^{\top}\nabla r(\beta_{0,k}^{*},\beta_{k}^{*})=0\right\} ,\label{Eqs: critical cone}
\end{equation}
and let
\[
\zeta_{k,n}=\sqrt{n}\left(\nabla R_{n}(\beta_{0,k}^{*},\beta_{k}^{*})-\nabla r(\beta_{0,k}^{*},\beta_{k}^{*})\right),\qquad Z_{k}\sim N(0,\Sigma_{k}),
\]
where $\Sigma_{k}=\mathrm{Var}\left(\nabla\ell(X,Y;\beta_{0,k}^{*},\beta_{k}^{*})\right)$
and $\zeta_{k,n}\rightsquigarrow Z_{k}$.

Define the function $\varphi_{k}:\mathbb{R}^{d+1}\to\mathbb{R}$ by
\begin{equation}
\varphi_{k}(\zeta)=\inf_{d\in \mathcal{C}_{k}(\beta_{0,k}^{*},\beta_{k}^{*})}\left\{ 2d^{\top}\zeta+d^{\top}\nabla^{2}r(\beta_{0,k}^{*},\beta_{k}^{*})d+\inf_{w\in \mathcal{T}_{\mathbb{T}_{k}}^{2}\left((\beta_{0,k}^{*},\beta_{k}^{*}),d\right)}w^{\top}\nabla r(\beta_{0,k}^{*},\beta_{k}^{*})\right\} ,\label{Eqs: phi k}
\end{equation}
and, when the minimiser is unique, let $\bar{d}_{k}(\zeta)$ denote
the unique minimiser of (\ref{Eqs: phi k}) with respect to $d$.

The displayed moment and positive-definiteness assumptions imply $\zeta_{k,n}\rightsquigarrow Z_k$ and $\nabla^2r(\theta_k^*)=H$. The theorem below retains uniqueness of the minimiser in (\ref{Eqs: phi k}) as an explicit assumption for the general curved constraints. In the Lasso case, $\mathbb T_k$ is polyhedral, the second-order tangent contribution vanishes, and the positive-definite quadratic term makes the resulting programme over the closed convex critical cone strictly convex; uniqueness then follows automatically.
For each $n$, let $(\widehat\beta_{0,k,n},\widehat\beta_{k,n})$ be a measurable selection from the empirical constrained argmin, using a fixed deterministic tie rule if the finite-sample minimiser is not unique, and let $\widehat v_{k,n}$ be its empirical risk value.

\begin{theorem}\label{Thm: Shapiro fixed k}
Fix $k$ and assume the conditions stated above. Suppose that for
every $\zeta\in\mathbb{R}^{d+1}$ the minimisation problem in (\ref{Eqs: phi k})
has a unique optimal solution $\bar{d}_{k}(\zeta)$. Then
\begin{align}
n\left[\hat{v}_{k,n}-R_{n}(\beta_{0,k}^{*},\beta_{k}^{*})\right] & \rightsquigarrow\frac{1}{2}\varphi_{k}(Z_{k}),\label{Eqs: value limit}\\
\hat{v}_{k,n} & =R_{n}(\beta_{0,k}^{*},\beta_{k}^{*})+\frac{1}{2n}\varphi_{k}(\zeta_{k,n})+o_{\mathrm{P}}(n^{-1}),\label{Eqs: value expansion}\\
\sqrt{n}\left(\begin{pmatrix}\hat{\beta}_{0,k,n}\\
\hat{\beta}_{k,n}
\end{pmatrix}-\begin{pmatrix}\beta_{0,k}^{*}\\
\beta_{k}^{*}
\end{pmatrix}\right) & \rightsquigarrow\bar{d}_{k}(Z_{k}).\label{Eqs: solution limit}
\end{align}
\end{theorem}

\begin{remark}
If $(\beta_{0,k}^{*},\beta_{k}^{*})$ is an interior point of $\mathbb{T}_{k}$, then $\nabla r(\beta_{0,k}^{*},\beta_{k}^{*})=0$, so $\mathcal{C}_{k}(\beta_{0,k}^{*},\beta_{k}^{*})=\mathcal{T}_{\mathbb{T}_{k}}(\beta_{0,k}^{*},\beta_{k}^{*})=\mathbb{R}^{d+1}$. In this case the second-order tangent term vanishes and $\bar{d}_{k}(Z_{k})=-\left\{ \nabla^{2}r(\beta_{0,k}^{*},\beta_{k}^{*})\right\}^{-1}Z_{k}$, yielding the usual asymptotically normal OLS-type limit distribution.

In the Lasso case, $\mathbb{T}_{k}$ is polyhedral, hence the critical cone $\mathcal{C}_{k}(\beta_{0,k}^{*},\beta_{k}^{*})$ is polyhedral. The limit in (\ref{Eqs: solution limit}) is then given by the unique minimiser of a strictly convex random quadratic programme over a polyhedral cone. Its distribution can assign positive probability to lower-dimensional faces of that cone.
\end{remark}



\begin{bluechange}
The scalar boundary case has a closed form, and the fixed-threshold law can be approximated by subsampling. Consider i.i.d. observations from the least-squares model without an intercept,
\[
Y=\beta^\circ X+\epsilon,\qquad \mathbb E(X\epsilon)=0,
\]
where $\beta^\circ>0$ is the population slope. Let $C>0$ be the constraint radius in $|\beta|\le C$, and consider the binding case $C=\beta^\circ$. Assume $0<q=\mathbb E(X^2)<\infty$, $\mathbb E(\epsilon^2)<\infty$, and $0<\tau^2=\operatorname{Var}(X\epsilon)<\infty$.
\begin{proposition}\label{Prop: scalar boundary}
Let $\widehat\beta_n$ be the constrained least-squares estimator under $|\beta|\le C$, where $C=\beta^\circ$. Then
\[
\sqrt n(\widehat\beta_n-\beta^\circ)\rightsquigarrow\min\{G,0\},
\qquad G\sim N\left(0,\frac{\tau^2}{q^2}\right).
\]
The limit has an atom of probability $1/2$ at zero and, on the negative half-line, the density of $N(0,\tau^2/q^2)$ without renormalisation. If $Z\sim N(0,4\tau^2)$ and $h=2q$, the corresponding value limit is
\[
n\{R_n(\widehat\beta_n)-R_n(\beta^\circ)\}\rightsquigarrow
\frac12\varphi(Z)=-\frac{Z^2}{2h}\,\mathds{1}\{Z>0\}.
\]
\end{proposition}

Write $\widehat\theta_{k,n}=(\widehat\beta_{0,k,n},\widehat\beta_{k,n}^\top)^\top$ for a measurable selection from the empirical constrained argmin. For a positive integer $b\le n$ and a subset $J\subset[n]$ of size $b$, let $\widehat\theta_{k,b,J}$ be the corresponding measurable selection, using the same deterministic tie rule, computed from the observations indexed by $J$.

\begin{proposition}\label{Prop: subsampling fixed k}
Fix $k$ and suppose that Theorem \ref{Thm: Shapiro fixed k} applies to this fixed candidate. Let $a\in\mathbb R^{d+1}$ be fixed and let $F_{k,a}$ denote the nondegenerate distribution function of $a^\top\bar d_k(Z_k)$. Let the deterministic integers $1\le b_n\le n$ satisfy $b_n\to\infty$ and $b_n/n\to0$, let $\mathcal J_{n,b_n}=\{J\subset[n]:|J|=b_n\}$, and define
\[
\widehat L_{n,b_n,a}(x)
=\binom{n}{b_n}^{-1}\sum_{J\in\mathcal J_{n,b_n}}
\mathds 1\!\left\{\sqrt{b_n}\,a^\top(\widehat\theta_{k,b_n,J}-\widehat\theta_{k,n})\le x\right\}.
\]
Then, at every continuity point $x$ of $F_{k,a}$,
\[
\widehat L_{n,b_n,a}(x)\xlongrightarrow{\mathbb P}F_{k,a}(x).
\]
Let $B_n\to\infty$ be deterministic. Conditional on the data, draw $B_n$ independent uniform subsets from $\mathcal J_{n,b_n}$ with replacement, or a uniform sample of $B_n\le\binom n{b_n}$ distinct subsets without replacement. The empirical distribution based on these subsets has the same pointwise limit.
\end{proposition}

The proposition applies \citet[Theorem 2.2.1 and Corollary 2.4.1]{politis1999subsampling} with normalising sequence $\tau_n=\sqrt n$; the condition $b_n/n\to0$ gives $\tau_{b_n}/\tau_n=\sqrt{b_n/n}\to0$. For Proposition \ref{Prop: scalar boundary}, the subsampling distribution therefore converges at every $x\ne0$ to the law of $\min\{G,0\}$. Lower-tail quantiles strictly below zero are consistently estimated at probability levels where the limiting distribution is continuous and strictly increasing, but no conclusion is asserted at the atom at zero.

For vector-valued uncertainty, \citet[Theorem 2.5.2]{politis1999subsampling} gives the corresponding norm construction when the limiting norm distribution is continuous at the quantile of interest. One admissible schedule is $b_n=\lfloor n^\gamma\rfloor$, $0<\gamma<1$; finite-sample sensitivity to $b_n$ should nevertheless be assessed.
\end{bluechange}
For selectors over a finite grid, exact selection yields the following consequence of Theorem \ref{Thm: Shapiro fixed k}. The interval-indexed case requires additional localisation control, discussed below.

\begin{corollary}\label{Cor: Shapiro selected} Assume that the conclusions
of Theorem \ref{Thm: Shapiro fixed k} hold at $k=k^{*}$, with unique
minimiser map $\bar{d}_{k^{*}}(\cdot)$. Let $\hat{K}_{n}$ take values
in $\left[m\right]$ and satisfy $\hat{K}_{n}\to k^{*}$ in probability, and
set $(\hat{\beta}_{0,n}^{\mathrm{sel}},\hat{\beta}_{n}^{\mathrm{sel}})=(\hat{\beta}_{0,\hat{K}_{n},n},\hat{\beta}_{\hat{K}_{n},n})$.
Then
\[
\sqrt{n}\left(\begin{pmatrix}\hat{\beta}_{0,n}^{\mathrm{sel}}\\
\hat{\beta}_{n}^{\mathrm{sel}}
\end{pmatrix}-\begin{pmatrix}\beta_{0,k^{*}}^{*}\\
\beta_{k^{*}}^{*}
\end{pmatrix}\right)\rightsquigarrow\bar{d}_{k^{*}}(Z_{k^{*}}).
\]
\end{corollary}
\begin{proof}
Since $\hat{K}_{n}$ takes values in $\left[m\right]$ and $\hat{K}_{n}\to k^{*}$
in probability, it follows that $\mathbb{P}(\hat{K}_{n}=k^{*})\to1$.
On the event $\{\hat{K}_{n}=k^{*}\}$ we have $(\hat{\beta}_{0,n}^{\mathrm{sel}},\hat{\beta}_{n}^{\mathrm{sel}})=(\hat{\beta}_{0,k^{*},n},\hat{\beta}_{k^{*},n})$,
hence
\[
\sqrt{n}\left(\begin{pmatrix}\hat{\beta}_{0,n}^{\mathrm{sel}}\\
\hat{\beta}_{n}^{\mathrm{sel}}
\end{pmatrix}-\begin{pmatrix}\beta_{0,k^{*}}^{*}\\
\beta_{k^{*}}^{*}
\end{pmatrix}\right)=\sqrt{n}\left(\begin{pmatrix}\hat{\beta}_{0,k^{*},n}\\
\hat{\beta}_{k^{*},n}
\end{pmatrix}-\begin{pmatrix}\beta_{0,k^{*}}^{*}\\
\beta_{k^{*}}^{*}
\end{pmatrix}\right)
\]
on that event. Combining with (\ref{Eqs: solution limit}) at $k=k^{*}$
and Slutsky's theorem yields the claim.
\end{proof}
\begin{bluechange}
In the interval
formulation $\hat{K}_{n}:\Omega\to[a,b]$, the consistency statement $\hat{K}_{n}\to k^{*}$
in probability does not by itself justify replacing $k$ by $\hat{K}_{n}$
in the second-order limit law of Theorem \ref{Thm: Shapiro fixed k}.
A standard sufficient substitution argument would require, for example, the localisation
$\hat{K}_{n}-k^{*}=o_{\mathrm{P}}(n^{-1/2})$, suitable local regularity of
$k\mapsto(\beta_{0,k}^{*},\beta_{k}^{*})$, and stochastic equicontinuity of the constrained estimator process near $k^*$. None of these properties is implied by first-order consistency alone. In constrained least squares, the estimator path can change its local form and the empirical value function can change curvature near the data-dependent threshold at which the constraint becomes inactive, so a sharper localisation rate requires a separate local analysis. Deriving the joint second-order limit of $\left(\hat{K}_{n},\hat{\beta}_{0,\hat{K}_{n},n},\hat{\beta}_{\hat{K}_{n},n}\right)$ in the interval-indexed setting is therefore beyond the scope of our framework.

\end{bluechange}

\section{Simulation Study} \label{Sect: Simulation}

\begin{bluechange}
The assessed methods are three implementable selectors: PanIC with the repeated two-way cross-signed multiplier rule in Appendix \ref{App: kappa calibration}, denoted PanIC-CF; the MSE-scale BIC-like criterion for Gaussian regression, with an analogue based on counts of non-zero slope parameters reported only as an exploratory comparator for logistic and Poisson regression; and ordinary minimum-loss five-fold cross-validation, denoted CV \citep{stone_cross-validatory_1974,hastie_elements_2009}. Proposition \ref{Prop: uniform kappa} shows that restricting the data-dependent multiplier to a fixed positive compact range preserves the first-order PanIC target. It does not assert that radius matching is optimal for prediction or support recovery, so these aspects are evaluated separately below. Support-informed choices of $\kappa$ are not used in the primary tables. The computational simulation specifications are detailed in Appendix \ref{App: reproducibility}.

\subsection{Simulation design}\label{Sect: Linear regression}
Every setting has covariates $X\sim N(0,\Sigma)$ in dimension $d=20$, with non-zero slope parameters indexed by $A^*=\{1,3,\ldots,19\}$. Define $\widetilde\beta_j=(-1)^{(j-1)/2}$ for $j\in A^*$ and $\widetilde\beta_j=0$ otherwise, and set
\[
\beta^*=\frac{\widetilde\beta}{\sqrt{\widetilde\beta^\top\Sigma\widetilde\beta}}.
\]
Thus the ten non-zero entries have equal magnitude and alternating signs, and $\beta^{*\top}\Sigma\beta^*=1$. Throughout the simulation study, $C^\star=\|\beta^*\|_1$ denotes the corresponding true constraint radius. For the independent design, $\Sigma=I_{20}$ and $C^\star=\sqrt{10}\approx3.1623$. For the correlated design, $\Sigma_{rs}=0.5^{|r-s|}$, whose diagonal entries are all one, and the rescaling gives $C^\star\approx3.9778$. Thus every simulated coordinate is already on the same population scale, and no empirical predictor standardisation is applied. Gaussian responses have intercept zero and noise standard deviation $\sigma=1$; their selection loss is divided by $\sigma^2$. Logistic responses have intercept zero. The Poisson design is independent Gaussian with
\[
\beta_0^*=\log(2)-\tfrac12\beta^{*\top}\Sigma\beta^*\approx0.1931,
\]
so that the marginal count mean is two. Conditional on $X$, the responses follow
\[
Y\mid X\sim
\begin{cases}
N(\beta_0^*+X^\top\beta^*,\sigma^2),&\text{Gaussian},\\
\operatorname{Bernoulli}\!\left(\{1+\exp[-(\beta_0^*+X^\top\beta^*)]\}^{-1}\right),&\text{logistic},\\
\operatorname{Poisson}\!\left(\exp(\beta_0^*+X^\top\beta^*)\right),&\text{Poisson}.
\end{cases}
\]
Each of the independent Gaussian, independent logistic, correlated Gaussian, correlated logistic, and independent Poisson designs is evaluated at $n\in\{500,1000\}$.

Each setting has $N=1{,}000$ Monte Carlo replications and an independent test sample of size 2,000 per replication. All selectors use the same fixed candidate grid $C_k=(k-1)/6$, $k=1,\ldots,121$, spanning $[0,20]$. This interval contains every simulated true radius. We use $\psi(C)=C/20$ and the PanIC shape $(1+\psi(C))\sqrt{\log n/n}$. The multiplier grid consists of 31 log-spaced values from $10^{-2}$ to $10^2$. The Gaussian BIC-like criterion uses $\kappa_D=1$, the cumulative exact non-zero parameter count, and $\varepsilon=10^{-6}$ in (\ref{Eqs: Modified BIC}). For compact table labels, the logistic and Poisson analogues are also denoted BIC-like. They use the likelihood-scale coefficient $1/2$ multiplying the same cumulative-count expression and are exploratory comparators outside the scope of Proposition \ref{Prop: BIC-like consistency}.


\subsection{Data-driven calibration and evaluation metrics}\label{Sect: calibration}
Within each replication, five balanced half-splits are generated from the recorded calibration seed and both directions are scored, giving ten calibration rows. Each row uses an ordinary training-half GLM to estimate coefficient signs and an independent ordinary validation-half GLM to estimate signed coefficient magnitudes; their cross-signed target is projected onto $[0,20]$ for scoring, while its raw value is retained for endpoint diagnostics. No regularised pilot is substituted. For the prespecified multiplier grid, we calculate the asymmetric discrepancy in (\ref{Eqs: fold radius discrepancy}), apply the largest-multiplier one-standard-error rule in (\ref{Eqs: kappa one standard error}), and refit on the full sample.

For the cross-validation selector, let $\mathcal I_1,\ldots,\mathcal I_V$, with $V=5$, be the balanced folds generated from the recorded CV seed. For candidate radius $C_k$, let $\widehat\vartheta_{-v,k}$ denote the constrained fit obtained without fold $\mathcal I_v$, and define
\[
\widehat L_{v,k}
=\frac{1}{|\mathcal I_v|}
\sum_{i\in\mathcal I_v}
\ell\{(X_i,Y_i);\widehat\vartheta_{-v,k}\}.
\]
Using the same family-specific selection loss as for the corresponding empirical risks, set
\[
\overline L_k=\frac{1}{V}\sum_{v=1}^V\widehat L_{v,k},
\qquad
\widehat k_{\rm CV}
=\min\left\{\argmin_{k\in[m]}\overline L_k\right\}.
\]
The ideal CV rule therefore chooses the smallest-radius minimiser in an exact tie. The numerical implementation treats criterion values within $10^{-12}$ of the minimum as tied. Its selected radius is refitted on the full sample before evaluation on the independent test sample, which does not enter the selection rule.

For the selected index set of non-zero slope parameters $\widehat A=\{j:\widehat\beta_j\ne0\}$, we report
\[
\mathrm{FP}=|\widehat A\setminus A^*|,\qquad
\mathrm{FN}=|A^*\setminus\widehat A|,
\]
\[
\mathrm{FPR}=\frac{\mathrm{FP}}{d-|A^*|},\qquad
\mathrm{FNR}=\frac{\mathrm{FN}}{|A^*|},\qquad
\mathrm{Exact}=\mathds{1}\{\widehat A=A^*\}.
\]
The zero/non-zero classification is applied to the fitted path coefficient after radius interpolation, without an additional magnitude threshold. The signed attained-radius error is $\|\widehat\beta\|_1-\|\beta^*\|_1$; under each correctly specified simulation, the data-generating slope $\beta^*$ is also the unregularised population minimiser $\beta^\circ$. Table \ref{Table: revised primary performance} reports this error and the selected grid radius $C_{\rm sel}$; the attained radius is therefore recoverable as the error plus $C^\star$. The two radii can differ when the constraint is inactive. Predictive performance is evaluated only on the independent test sample. For Gaussian regression, test deviance is the mean squared error divided by $\sigma^2$. For logistic regression, with fitted probabilities $\widehat p_i$, it is
\[
\frac{2}{n_{\rm test}}\sum_{i=1}^{n_{\rm test}}
\left[
Y_i\log\left(\frac{Y_i}{\widehat p_i}\right)
+(1-Y_i)\log\left\{\frac{1-Y_i}{1-\widehat p_i}\right\}
\right],
\]
and for Poisson regression, with fitted means $\widehat\mu_i$, it is
\[
\frac{2}{n_{\rm test}}\sum_{i=1}^{n_{\rm test}}
\left[
Y_i\log\left(\frac{Y_i}{\widehat\mu_i}\right)-(Y_i-\widehat\mu_i)
\right].
\]
We use the convention $0\log0=0$. No support or test-sample quantity enters the PanIC-CF calibration. Parenthetical uncertainty is the Monte Carlo standard error (MCSE), computed as the sample standard deviation across independent replications divided by the square root of the number of available replications; paired contrasts are formed within replication before this calculation.

\subsection{Primary results}\label{Sect: Discussion}
Tables \ref{Table: revised primary support} and
\ref{Table: revised primary performance} report the support, radius, and
independent-test summaries for the assessed methods.
% BEGIN PANIC INLINE TABLE: table_primary_support.tex
\begin{table}[H]
\centering
\caption{Support recovery for the assessed methods over 1,000 replications; parentheses give Monte Carlo standard errors. FPR and FNR use the ten inactive and ten active slopes, respectively.}
\label{Table: revised primary support}
\begin{adjustbox}{width=\textwidth}
\scriptsize
\begin{tabular}{llrrrrrr}
\hline
Setting & Method & FP & FN & FPR & FNR & Exact & Total support error \\
\hline
Gaussian independent, $n=500$ & PanIC-CF & 5.99 (0.06) & 0.00 (0.00) & 0.599 (0.006) & 0.000 (0.000) & 0.002 (0.001) & 5.99 (0.06) \\
 & BIC-like & 2.68 (0.05) & 0.00 (0.00) & 0.268 (0.005) & 0.000 (0.000) & 0.076 (0.008) & 2.68 (0.05) \\
 & CV & 6.65 (0.06) & 0.00 (0.00) & 0.665 (0.006) & 0.000 (0.000) & 0.000 (0.000) & 6.65 (0.06) \\
Gaussian independent, $n=1000$ & PanIC-CF & 6.07 (0.06) & 0.00 (0.00) & 0.607 (0.006) & 0.000 (0.000) & 0.001 (0.001) & 6.07 (0.06) \\
 & BIC-like & 2.34 (0.05) & 0.00 (0.00) & 0.234 (0.005) & 0.000 (0.000) & 0.126 (0.010) & 2.34 (0.05) \\
 & CV & 6.64 (0.06) & 0.00 (0.00) & 0.664 (0.006) & 0.000 (0.000) & 0.000 (0.000) & 6.64 (0.06) \\
Logistic independent, $n=500$ & PanIC-CF & 6.27 (0.06) & 0.05 (0.01) & 0.627 (0.006) & 0.005 (0.001) & 0.003 (0.002) & 6.32 (0.06) \\
 & BIC-like & 1.85 (0.06) & 1.43 (0.09) & 0.185 (0.006) & 0.143 (0.009) & 0.075 (0.008) & 3.27 (0.08) \\
 & CV & 6.35 (0.06) & 0.05 (0.01) & 0.635 (0.006) & 0.005 (0.001) & 0.000 (0.000) & 6.40 (0.06) \\
Logistic independent, $n=1000$ & PanIC-CF & 6.53 (0.06) & 0.00 (0.00) & 0.653 (0.006) & 0.000 (0.000) & 0.001 (0.001) & 6.53 (0.06) \\
 & BIC-like & 2.37 (0.06) & 0.01 (0.00) & 0.237 (0.006) & 0.001 (0.000) & 0.128 (0.011) & 2.37 (0.05) \\
 & CV & 6.64 (0.06) & 0.00 (0.00) & 0.664 (0.006) & 0.000 (0.000) & 0.000 (0.000) & 6.64 (0.06) \\
Gaussian AR(1), $n=500$ & PanIC-CF & 5.91 (0.06) & 0.00 (0.00) & 0.591 (0.006) & 0.000 (0.000) & 0.006 (0.002) & 5.91 (0.06) \\
 & BIC-like & 2.64 (0.05) & 0.00 (0.00) & 0.264 (0.005) & 0.000 (0.000) & 0.078 (0.008) & 2.64 (0.05) \\
 & CV & 6.61 (0.06) & 0.00 (0.00) & 0.661 (0.006) & 0.000 (0.000) & 0.002 (0.001) & 6.61 (0.06) \\
Gaussian AR(1), $n=1000$ & PanIC-CF & 5.92 (0.06) & 0.00 (0.00) & 0.593 (0.006) & 0.000 (0.000) & 0.001 (0.001) & 5.92 (0.06) \\
 & BIC-like & 2.47 (0.05) & 0.00 (0.00) & 0.247 (0.005) & 0.000 (0.000) & 0.111 (0.010) & 2.47 (0.05) \\
 & CV & 6.55 (0.07) & 0.00 (0.00) & 0.655 (0.007) & 0.000 (0.000) & 0.001 (0.001) & 6.55 (0.07) \\
Logistic AR(1), $n=500$ & PanIC-CF & 6.25 (0.06) & 0.03 (0.01) & 0.625 (0.006) & 0.003 (0.001) & 0.000 (0.000) & 6.27 (0.06) \\
 & BIC-like & 1.97 (0.06) & 1.63 (0.10) & 0.197 (0.006) & 0.163 (0.010) & 0.067 (0.008) & 3.60 (0.09) \\
 & CV & 6.41 (0.06) & 0.02 (0.00) & 0.641 (0.006) & 0.002 (0.000) & 0.000 (0.000) & 6.43 (0.06) \\
Logistic AR(1), $n=1000$ & PanIC-CF & 6.47 (0.06) & 0.00 (0.00) & 0.647 (0.006) & 0.000 (0.000) & 0.001 (0.001) & 6.47 (0.06) \\
 & BIC-like & 2.35 (0.05) & 0.00 (0.00) & 0.235 (0.005) & 0.000 (0.000) & 0.123 (0.010) & 2.36 (0.05) \\
 & CV & 6.62 (0.06) & 0.00 (0.00) & 0.662 (0.006) & 0.000 (0.000) & 0.000 (0.000) & 6.62 (0.06) \\
Poisson independent, $n=500$ & PanIC-CF & 6.27 (0.06) & 0.00 (0.00) & 0.627 (0.006) & 0.000 (0.000) & 0.001 (0.001) & 6.27 (0.06) \\
 & BIC-like & 2.85 (0.06) & 0.00 (0.00) & 0.285 (0.006) & 0.000 (0.000) & 0.074 (0.008) & 2.85 (0.06) \\
 & CV & 6.74 (0.06) & 0.00 (0.00) & 0.674 (0.006) & 0.000 (0.000) & 0.000 (0.000) & 6.74 (0.06) \\
Poisson independent, $n=1000$ & PanIC-CF & 6.37 (0.06) & 0.00 (0.00) & 0.637 (0.006) & 0.000 (0.000) & 0.001 (0.001) & 6.37 (0.06) \\
 & BIC-like & 2.69 (0.06) & 0.00 (0.00) & 0.269 (0.006) & 0.000 (0.000) & 0.107 (0.010) & 2.69 (0.06) \\
 & CV & 6.90 (0.07) & 0.00 (0.00) & 0.690 (0.007) & 0.000 (0.000) & 0.000 (0.000) & 6.90 (0.07) \\
\hline
\end{tabular}
\end{adjustbox}
\end{table}
% END PANIC INLINE TABLE: table_primary_support.tex


% BEGIN PANIC INLINE TABLE: table_primary_performance.tex
\begin{table}[H]
\centering
\caption{Selected radii and independent-test performance for the assessed methods over 1,000 replications; parentheses give Monte Carlo standard errors.}
\label{Table: revised primary performance}
\begin{adjustbox}{width=\textwidth}
\scriptsize
\begin{tabular}{llrrr}
\hline
Setting & Method & Signed attained-radius error & $C_{\rm sel}$ & Test deviance \\
\hline
Gaussian independent, $n=500$ & PanIC-CF & -0.067 (0.007) & 3.095 (0.007) & 1.0396 (0.0012) \\
 & BIC-like & -0.366 (0.008) & 2.796 (0.008) & 1.0514 (0.0013) \\
 & CV & 0.014 (0.007) & 3.176 (0.007) & 1.0393 (0.0012) \\
Gaussian independent, $n=1000$ & PanIC-CF & -0.049 (0.005) & 3.113 (0.005) & 1.0193 (0.0011) \\
 & BIC-like & -0.301 (0.005) & 2.861 (0.005) & 1.0269 (0.0011) \\
 & CV & 0.002 (0.005) & 3.164 (0.005) & 1.0193 (0.0011) \\
Logistic independent, $n=500$ & PanIC-CF & -0.134 (0.017) & 3.028 (0.017) & 1.2356 (0.0007) \\
 & BIC-like & -1.208 (0.029) & 1.954 (0.029) & 1.2620 (0.0015) \\
 & CV & -0.104 (0.017) & 3.059 (0.017) & 1.2355 (0.0007) \\
Logistic independent, $n=1000$ & PanIC-CF & -0.091 (0.011) & 3.071 (0.011) & 1.2177 (0.0005) \\
 & BIC-like & -0.748 (0.014) & 2.414 (0.014) & 1.2251 (0.0005) \\
 & CV & -0.049 (0.012) & 3.113 (0.012) & 1.2177 (0.0006) \\
Gaussian AR(1), $n=500$ & PanIC-CF & -0.088 (0.008) & 3.889 (0.008) & 1.0373 (0.0011) \\
 & BIC-like & -0.456 (0.010) & 3.522 (0.010) & 1.0488 (0.0012) \\
 & CV & 0.020 (0.009) & 3.998 (0.009) & 1.0371 (0.0011) \\
Gaussian AR(1), $n=1000$ & PanIC-CF & -0.064 (0.006) & 3.914 (0.006) & 1.0176 (0.0011) \\
 & BIC-like & -0.355 (0.007) & 3.623 (0.007) & 1.0244 (0.0011) \\
 & CV & 0.007 (0.006) & 3.984 (0.006) & 1.0175 (0.0011) \\
Logistic AR(1), $n=500$ & PanIC-CF & -0.208 (0.022) & 3.770 (0.022) & 1.2370 (0.0007) \\
 & BIC-like & -1.551 (0.040) & 2.427 (0.040) & 1.2670 (0.0017) \\
 & CV & -0.122 (0.021) & 3.856 (0.021) & 1.2366 (0.0007) \\
Logistic AR(1), $n=1000$ & PanIC-CF & -0.115 (0.014) & 3.862 (0.014) & 1.2173 (0.0006) \\
 & BIC-like & -0.918 (0.017) & 3.059 (0.017) & 1.2241 (0.0005) \\
 & CV & -0.064 (0.015) & 3.914 (0.015) & 1.2174 (0.0006) \\
Poisson independent, $n=500$ & PanIC-CF & -0.040 (0.005) & 3.123 (0.005) & 1.0958 (0.0011) \\
 & BIC-like & -0.278 (0.006) & 2.884 (0.006) & 1.1107 (0.0013) \\
 & CV & 0.003 (0.005) & 3.165 (0.005) & 1.0956 (0.0011) \\
Poisson independent, $n=1000$ & PanIC-CF & -0.022 (0.004) & 3.140 (0.004) & 1.0739 (0.0009) \\
 & BIC-like & -0.211 (0.004) & 2.952 (0.004) & 1.0819 (0.0010) \\
 & CV & 0.010 (0.004) & 3.172 (0.004) & 1.0738 (0.0009) \\
\hline
\end{tabular}
\end{adjustbox}
\end{table}
% END PANIC INLINE TABLE: table_primary_performance.tex


PanIC-CF had lower mean total support error than CV in all ten settings, mainly through fewer false positives; the paired mean differences (PanIC-CF minus CV) ranged from $-0.697$ to $-0.086$. The absolute difference between their mean test deviances was at most $0.0004$.


The rows labelled BIC-like had mean total support errors from $2.34$ to $3.60$ and exact-recovery proportions from $0.067$ to $0.128$, together with larger mean test deviance than CV. These comparisons reflect different support--prediction trade-offs and do not establish a universal ranking: matching the unregularised population radius, recovering a sparse support, and minimising predictive loss are distinct objectives.
\end{bluechange}

\begin{bluechange}
\subsection{Calibration and numerical diagnostics}
Table \ref{Table: revised calibration} reports the calibration targets and
selected multipliers.
% BEGIN PANIC INLINE TABLE: table_calibration.tex
\begin{table}[H]
\centering
\caption{PanIC-CF calibration over $N=1{,}000$ replications. Raw cross-signed targets are first averaged within each data-set replication; parentheses give Monte Carlo standard errors across replications. Raw bias is the mean target minus $C^\star$; multiplier entries are median [Q1,Q3].}
\label{Table: revised calibration}
\begin{adjustbox}{width=\textwidth}
\scriptsize
\begin{tabular}{lrrrl}
\hline
Setting & $C^\star$ & Mean raw target & Raw bias & $\widehat\kappa$ [Q1,Q3] \\
\hline
Gaussian independent, $n=500$ & 3.162 & 3.170 (0.006) & 0.007 (0.006) & 8.577 [6.310, 8.577] \\
Gaussian independent, $n=1000$ & 3.162 & 3.158 (0.004) & -0.005 (0.004) & 8.577 [6.310, 8.577] \\
Logistic independent, $n=500$ & 3.162 & 3.433 (0.019) & 0.271 (0.019) & 1.848 [1.359, 1.848] \\
Logistic independent, $n=1000$ & 3.162 & 3.307 (0.011) & 0.145 (0.011) & 1.359 [1.359, 1.848] \\
Gaussian AR(1), $n=500$ & 3.978 & 3.987 (0.008) & 0.009 (0.008) & 6.310 [6.310, 8.577] \\
Gaussian AR(1), $n=1000$ & 3.978 & 3.975 (0.006) & -0.003 (0.006) & 6.310 [4.642, 8.577] \\
Logistic AR(1), $n=500$ & 3.978 & 4.282 (0.024) & 0.305 (0.024) & 1.359 [1.000, 1.848] \\
Logistic AR(1), $n=1000$ & 3.978 & 4.174 (0.015) & 0.197 (0.015) & 1.359 [1.000, 1.359] \\
Poisson independent, $n=500$ & 3.162 & 3.171 (0.005) & 0.008 (0.005) & 4.642 [4.642, 6.310] \\
Poisson independent, $n=1000$ & 3.162 & 3.164 (0.003) & 0.001 (0.003) & 4.642 [4.642, 6.310] \\
\hline
\end{tabular}
\end{adjustbox}
\end{table}
% END PANIC INLINE TABLE: table_calibration.tex


All 10,000 primary replications completed without a fitting or calibration failure, solver warning, multiplier default, or nontrivial projection of a raw cross-signed target. PanIC-CF selected neither a multiplier nor a radius endpoint.



\subsection{Grid sensitivity}
% BEGIN PANIC INLINE TABLE: table_grid_sensitivity.tex
\begin{table}[H]
\centering
\caption{Radius-grid sensitivity for the assessed methods in the independent Gaussian setting with $n=1000$ and $N=500$ common-random-number replications. Each grid spans $[0,20]$; entries are means with Monte Carlo standard errors in parentheses.}
\label{Table: grid sensitivity}
\begin{adjustbox}{width=\textwidth}
\scriptsize
\begin{tabular}{rrlrrrr}
\hline
$m$ & Spacing & Method & Total support error & Exact & Test deviance & Selected radius \\
\hline
61 & 0.333 & PanIC-CF & 6.47 (0.10) & 0.000 (0.000) & 1.019 (0.001) & 3.149 (0.008) \\
61 & 0.333 & BIC-like & 2.71 (0.08) & 0.124 (0.015) & 1.027 (0.001) & 2.867 (0.009) \\
61 & 0.333 & CV & 6.80 (0.10) & 0.000 (0.000) & 1.019 (0.001) & 3.183 (0.008) \\
121 & 0.167 & PanIC-CF & 6.24 (0.09) & 0.002 (0.002) & 1.019 (0.001) & 3.137 (0.007) \\
121 & 0.167 & BIC-like & 2.44 (0.07) & 0.094 (0.013) & 1.026 (0.001) & 2.875 (0.008) \\
121 & 0.167 & CV & 6.77 (0.09) & 0.000 (0.000) & 1.019 (0.001) & 3.183 (0.008) \\
241 & 0.083 & PanIC-CF & 6.10 (0.09) & 0.000 (0.000) & 1.019 (0.001) & 3.126 (0.007) \\
241 & 0.083 & BIC-like & 2.23 (0.07) & 0.124 (0.015) & 1.025 (0.001) & 2.881 (0.008) \\
241 & 0.083 & CV & 6.67 (0.09) & 0.000 (0.000) & 1.019 (0.001) & 3.182 (0.007) \\
\hline
\end{tabular}
\end{adjustbox}
\end{table}
% END PANIC INLINE TABLE: table_grid_sensitivity.tex


The grid-sensitivity study uses 500 replications of the Gaussian independent setting at $n=1000$, with common random numbers and 61, 121, and 241 equally spaced radii on $[0,20]$ (Table \ref{Table: grid sensitivity}). Across these grids, mean PanIC-CF test deviance remained between $1.0186$ and $1.0192$, while mean support error decreased from $6.468$ to $6.100$ and mean shared end-to-end per-replication runtime for all three methods increased from $0.0652$ to $0.1261$ seconds. The paired support-error changes were $-0.226$ (MCSE $0.086$) from 61 to 121 points and $-0.142$ ($0.061$) from 121 to 241 points. Thus prediction was stable over these three grids, while support error, selected radius, and computation retained some grid-resolution dependence. This finite comparison is not an asymptotic refining-grid claim.

\subsection{Boundary-law and subsampling illustration}
We illustrate the boundary limit in Proposition \ref{Prop: scalar boundary} and the scalar subsampling construction following Proposition \ref{Prop: subsampling fixed k} with $\beta^\circ=C=1$ and independent $X,\epsilon\sim N(0,1)$. The numerical illustration uses the upper-bound estimator $\widehat\beta_n=\min\{\widehat\beta_{\mathrm{OLS}},1\}$. It agrees with the estimator under $|\beta|\le1$ whenever $\widehat\beta_{\mathrm{OLS}}\ge-1$, an event whose probability tends to one, and therefore has the same boundary limit. Thus $q=\tau^2=1$ and the limiting continuous component is standard Gaussian on the negative half-line. For each of $n=500$ and $n=1000$, 40,000 independent samples were generated. Table \ref{Table: boundary illustration} and the first two panels of Figure \ref{fig:boundary subsampling} show that the empirical boundary mass and lower quantiles are close to their analytic values.

% BEGIN PANIC INLINE TABLE: table_boundary.tex
\begin{table}[H]
\centering
\caption{Scalar binding-boundary illustration with 40,000 replications. The limiting law has a point mass $1/2$ at zero and a standard-Gaussian density on the negative half-line. MCSE denotes the binomial Monte Carlo standard error of the atom estimate.}
\label{Table: boundary illustration}
\small
\begin{tabular}{rrrrrr}
\hline
$n$ & Atom & MCSE & 10th percentile & 25th percentile & Limit atom \\
\hline
500 & 0.5021 & 0.0025 & -1.282 & -0.665 & 0.5 \\
1000 & 0.4989 & 0.0025 & -1.294 & -0.681 & 0.5 \\
Limit & 0.5 & -- & -1.282 & -0.674 & 0.5 \\
\hline
\end{tabular}
\end{table}
% END PANIC INLINE TABLE: table_boundary.tex


For the subsampling illustration, a single full sample of size $n=1000$ was generated from the prespecified seed. Its unconstrained least-squares estimate was $1.0307$, so the upper-bound estimate was the binding value $\widehat\beta_n=1$. For each $b\in\{50,100,200\}$, $B_n=10{,}000$ subsets were sampled uniformly without replacement. Table \ref{Table: subsampling sensitivity} reports lower quantiles and CDF values at three negative continuity points. In this realised sample, $b/n=0.05$ gave the closest empirical CDF value at each displayed point. This finite comparison is consistent with, but is not implied monotonically by, the asymptotic condition $b/n\to0$.

% BEGIN PANIC INLINE TABLE: table_subsampling.tex
\begin{table}[H]
\centering
\caption{Subsampling sensitivity for one binding full sample of size $n=1000$, using 10,000 random subsamples at each $b$. The final three columns compare empirical cdf values with the limit cdf at continuity points.}
\label{Table: subsampling sensitivity}
\begin{adjustbox}{width=\textwidth}
\small
\begin{tabular}{rrrrrrrr}
\hline
$b$ & $b/n$ & Atom at zero & 10th pct. & 25th pct. & $F(-1)$ & $F(-0.5)$ & $F(-0.1)$ \\
\hline
50 & 0.05 & 0.578 & -1.104 & -0.476 & 0.118 & 0.243 & 0.385 \\
100 & 0.10 & 0.616 & -0.975 & -0.380 & 0.096 & 0.213 & 0.349 \\
200 & 0.20 & 0.673 & -0.801 & -0.217 & 0.070 & 0.169 & 0.291 \\
Limit & 0 & 0.500 & -1.282 & -0.674 & 0.159 & 0.309 & 0.460 \\
\hline
\end{tabular}
\end{adjustbox}
\end{table}
% END PANIC INLINE TABLE: table_subsampling.tex


\begin{figure}[H]
    \centering
    \safeincludegraphics[width=\textwidth]{figure_boundary_subsampling.pdf}
    \caption{Scalar binding-boundary and subsampling illustration \blue{using $\widehat\beta_n=\min\{\widehat\beta_{\mathrm{OLS}},1\}$}. The first two panels compare the empirical negative component and boundary atom of $\sqrt n(\widehat\beta_n-\beta^\circ)$ with the analytic limit for $n=500$ and $n=1000$. The third panel compares subsampling CDFs from the prespecified-seed full sample of size $n=1000$, using $B_n=10{,}000$ repeated random subsets for each $b\in\{50,100,200\}$, with the limiting CDF.}
    \label{fig:boundary subsampling}
\end{figure}

\subsection{Runtime}
\begin{figure}[H]
    \centering
    \safeincludegraphics[width=\textwidth]{figure_runtime.pdf}
    \caption{Elapsed fitting time for the assessed methods at $n=1000$ on the 121-point radius grid. Each boxplot contains 20 one-worker timed runs.}
    \label{fig:runtime}
\end{figure}

Each timed run used one worker after two untimed warm-up runs; data generation and independent-test evaluation were excluded. PanIC-CF had higher mean elapsed time than CV because it fits ten half-sample paths, the associated ordinary training and validation GLMs, and one full-sample path, whereas CV fits five fold-specific paths and one full-sample path. The PanIC-CF and CV means (MCSEs), respectively, were $0.0457$ ($0.0002$) and $0.0254$ ($0.0002$) seconds for Gaussian regression, $0.1607$ ($0.0022$) and $0.0970$ ($0.0006$) seconds for logistic regression, and $0.1534$ ($0.0009$) and $0.0861$ ($0.0006$) seconds for Poisson regression. The corresponding BIC-like means were $0.0046$ ($0.0001$), $0.0180$ ($0.0002$), and $0.0161$ ($0.0002$) seconds. These timings compare the stated implementations on one machine, not general computational complexity.
\end{bluechange}

\section{Conclusion}\label{Sect: Conclusion}
\begin{bluechange}
This paper develops information-criterion methods for selecting the constraint level in regularised generalised linear models. The PanIC regularity conditions are verified for linear, logistic, Poisson, and Gamma regression with an intercept constrained to a common compact interval, and the Gamma result permits joint estimation of a shape parameter on a fixed compact interval bounded away from zero. For regularised linear least squares, the paper also gives a monotone MSE-scale BIC-like construction whose consistency argument separates the generic complexity ordering from model-specific degrees-of-freedom interpretations.

The first-order theory covers deterministic continuum-indexed nested families and finite random candidate sets satisfying explicit one-sided coverage, empirical-risk, and penalty-separation conditions. At the second order, the constrained least-squares estimator is described by a random quadratic programme over a critical cone; the scalar binding-boundary case has a closed form, and subsampling consistently approximates nondegenerate scalar projections for a fixed candidate. Exact finite-grid selection transfers the fixed-candidate estimator law, whereas a continuum-selected second-order limit would require a separate localisation and stochastic-equicontinuity analysis.

The compact-multiplier result supports the repeated two-way cross-signed rule in Appendix \ref{App: kappa calibration}. In the ten-setting comparison against CV, PanIC-CF had lower mean support error in every setting, while the absolute differences in mean test deviance were at most $0.0004$. The BIC-like comparators exhibited a different support--prediction trade-off. These simulations do not establish a general dominance ordering because constraint-radius estimation, support recovery, and prediction are distinct targets. The binding-boundary and subsampling experiments illustrate the non-Gaussian limit and its approximation at continuity points. Further work includes a joint second-order limit for a continuum-selected threshold and its estimator, and high-dimensional calibration when ordinary unregularised half-sample fits are unavailable.
\end{bluechange}

\section*{Acknowledgements}
Funding was provided by Australian Research Council grants DP230100905 and DP250100860.

\appendix

\section[Repeated cross-signed calibration of the PanIC multiplier]{Repeated cross-signed calibration of the PanIC \mbox{multiplier}}\label{App: kappa calibration}
\begin{bluechange}
This appendix specifies the data-dependent multiplier rule used in Section \ref{Sect: calibration}. Each calibration row combines two ordinary unregularised estimators fitted on disjoint halves: the training-half estimator supplies coefficient signs and the validation-half estimator supplies signed coefficient magnitudes. The calibration discrepancy is centred on the smallest risk-optimal constraint radius; the finite multiplier search is a calibration device rather than a separate consistency theorem. The displays below use the compact theoretical parameter domains. Section \ref{Sect: Computations} describes the unconstrained numerical intercept and the localisation condition under which the theoretical and numerical intercept treatments agree asymptotically.

For linear, logistic, and Poisson regression, write $\vartheta=(\beta_0,\beta^\top)^\top$, where $\beta$ is the slope vector. When the Gamma shape is estimated jointly, write $\vartheta=(\beta_0,\beta^\top,\nu)^\top$, with $\nu\in\mathbb V$ treated as a nuisance parameter. In either case, the constraint functional below depends only on $\beta$. Let
\[
g(\beta)=
\begin{cases}
\|\beta\|_1, & \text{for Lasso},\\
\|\beta\|_2^2, & \text{for Ridge},\\
\alpha\|\beta\|_1+(1-\alpha)\|\beta\|_2^2, & \text{for Elastic-net}.
\end{cases}
\]
Let $\Theta_{\mathrm u}=\mathbb U\times\mathbb R^d$ for linear, logistic, and Poisson regression, and let $\Theta_{\mathrm u}=\mathbb U\times\mathbb R^d\times\mathbb V$ when the Gamma shape is estimated jointly. Thus the slopes are unconstrained, the intercept remains in the common compact interval $\mathbb U$ but is unpenalised, and any model-specific nuisance restriction is retained.

\subsection{Population radius and the unregularised pilot}
Assume that the population risk has a unique unregularised minimiser $\vartheta^\circ=\argmin_{\vartheta\in\Theta_{\mathrm u}}r(\vartheta)$. For linear, logistic, and Poisson regression, write $\vartheta^\circ=(\beta_0^\circ,\beta^{\circ\top})^\top$; when the Gamma shape is estimated jointly, write $\vartheta^\circ=(\beta_0^\circ,\beta^{\circ\top},\nu^\circ)^\top$, with $\nu^\circ\in\mathbb V$. In either case, let
\[
C^\circ=g(\beta^\circ),
\]
with $\beta_0^\circ$ in the interior of the common intercept interval $\mathbb U$. For $C\ge0$, let
\[
\mathbb T(C)=\{\vartheta\in\Theta_{\mathrm u}:g(\beta)\le C\},
\qquad
v(C)=\min_{\vartheta\in\mathbb T(C)}r(\vartheta),
\]
and assume that the displayed minimum exists on the constraint interval under consideration.

\begin{proposition}\label{Prop: population pilot radius}
Under the preceding assumptions,
\[
v(C)=r(\vartheta^\circ)\quad\text{if and only if}\quad C\ge C^\circ.
\]
Consequently, for a fixed grid $0\le C_1<\cdots<C_m$ with $m\ge2$ and $C_1\le C^\circ\le C_m$, the smallest risk-optimal candidate index is
\[
k_C^*=\min\{k\in[m]:C_k\ge C^\circ\}.
\]
\end{proposition}

\begin{proof}
If $C\ge C^\circ$, then $\vartheta^\circ\in\mathbb T(C)$, so $v(C)\le r(\vartheta^\circ)$. The reverse inequality follows because $\vartheta^\circ$ is the global minimiser over $\Theta_{\mathrm u}$. Hence equality holds. Conversely, suppose $C<C^\circ$ and $v(C)=r(\vartheta^\circ)$. Any minimiser attaining $v(C)$ is then another global minimiser of $r$ over $\Theta_{\mathrm u}$. Uniqueness forces it to equal $\vartheta^\circ$, which contradicts $g(\beta^\circ)=C^\circ>C$. The grid statement follows immediately.
\end{proof}

Here $Z_i=(X_i,Y_i)$. For a sample indexed by a nonempty set $S$, define
\[
R_S(\vartheta)=\frac{1}{|S|}\sum_{i\in S}\ell(Z_i;\vartheta),
\qquad
\widetilde\vartheta_S=\argmin_{\vartheta\in\Theta_{\mathrm u}}R_S(\vartheta),
\qquad
\widetilde C_S=g(\widetilde\beta_S),
\]
whenever the ordinary unregularised empirical-risk minimiser exists uniquely and is finite. For the following rate statements, assume that this event is measurable and has probability tending to one as $|S|\to\infty$, and use a measurable arbitrary completion on its complement. This completion is only a convention for stating a rate and does not substitute a regularised pilot. The following condition makes explicit what is required of the completed estimator:
\begin{equation}\label{Eqs: unregularised pilot rate}
\sqrt{|S|}\,(\widetilde\beta_S-\beta^\circ)=O_{\mathbb P}(1).
\end{equation}

\begin{proposition}\label{Prop: unregularised pilot radius}
If (\ref{Eqs: unregularised pilot rate}) holds, then
\[
\widetilde C_S-C^\circ=O_{\mathbb P}(|S|^{-1/2}).
\]
The conclusion holds for each of the Lasso, Ridge, and Elastic-net functionals above.
\end{proposition}

\begin{proof}
For Lasso,
\[
|g(\widetilde\beta_S)-g(\beta^\circ)|
\le \|\widetilde\beta_S-\beta^\circ\|_1
\le \sqrt d\,\|\widetilde\beta_S-\beta^\circ\|_2.
\]
For Ridge,
\[
\left|\|\widetilde\beta_S\|_2^2-\|\beta^\circ\|_2^2\right|
\le(\|\widetilde\beta_S\|_2+\|\beta^\circ\|_2)
\|\widetilde\beta_S-\beta^\circ\|_2,
\]
and $\|\widetilde\beta_S\|_2=O_{\mathbb P}(1)$ follows from (\ref{Eqs: unregularised pilot rate}). The Elastic-net bound is the corresponding convex combination of these two bounds. Multiplication by $\sqrt{|S|}$ proves the claim.
\end{proof}

For the Lasso functional, the plug-in radius $\|\widetilde\beta_S\|_1$ contains the absolute-value contribution $|\widetilde\beta_{S,j}|$ from every null coordinate. The implemented calibration instead separates sign estimation from coefficient estimation. For disjoint nonempty samples $S$ and $T$, define $\operatorname{sign}(0)=0$ and
\begin{equation}\label{Eqs: cross signed radius}
\widetilde C_{S,T}^{\times}
=\sum_{j=1}^d
\operatorname{sign}(\widetilde\beta_{S,j})\widetilde\beta_{T,j},
\qquad
\widetilde C_{S,T}^{\mathrm{proj}}
=\Pi_{[C_1,C_m]}(\widetilde C_{S,T}^{\times}),
\end{equation}
where $\Pi_{[C_1,C_m]}(x)=\min\{C_m,\max(C_1,x)\}$.

\begin{proposition}\label{Prop: cross signed pilot radius}
Suppose $d$ is fixed, $C^\circ=\|\beta^\circ\|_1\in[C_1,C_m]$, and the ordinary estimators on $S$ and $T$ satisfy
\[
\sqrt{|S|}(\widetilde\beta_S-\beta^\circ)=O_{\mathbb P}(1),
\qquad
\sqrt{|T|}(\widetilde\beta_T-\beta^\circ)=O_{\mathbb P}(1).
\]
If $|S|$ and $|T|$ are both of order $N$, then
\[
\widetilde C_{S,T}^{\times}-C^\circ=O_{\mathbb P}(N^{-1/2}),
\qquad
\widetilde C_{S,T}^{\mathrm{proj}}-C^\circ=O_{\mathbb P}(N^{-1/2}).
\]
This is a pointwise fixed-dimensional statement for a fixed $\beta^\circ$; it is not uniform over sequences whose nonzero coordinates approach zero.
\end{proposition}

\begin{proof}
Let $A=\{j:\beta_j^\circ\ne0\}$. Root-$|S|$ consistency and fixed non-zero coefficients imply
\[
\mathbb P\!\left\{
\operatorname{sign}(\widetilde\beta_{S,j})
=\operatorname{sign}(\beta_j^\circ)
\text{ for every }j\in A
\right\}\longrightarrow1.
\]
On this event,
\begin{align*}
\widetilde C_{S,T}^{\times}-C^\circ
&=\sum_{j\in A}\operatorname{sign}(\beta_j^\circ)
(\widetilde\beta_{T,j}-\beta_j^\circ)
+\sum_{j\notin A}\operatorname{sign}(\widetilde\beta_{S,j})
\widetilde\beta_{T,j}.
\end{align*}
Both sums are bounded in absolute value by a fixed multiple of
$\|\widetilde\beta_T-\beta^\circ\|_2$ and are therefore
$O_{\mathbb P}(N^{-1/2})$. The complement of the sign-consistency event has probability tending to zero, which proves the first claim. Projection onto an interval containing $C^\circ$ is nonexpansive, so
\[
|\widetilde C_{S,T}^{\mathrm{proj}}-C^\circ|
\le |\widetilde C_{S,T}^{\times}-C^\circ|,
\]
and the second claim follows.
\end{proof}

Using both directions of each split follows the general cross-fitting device of exchanging training and evaluation roles \citep{ChernozhukovEtAl2018DML}, while repeated splits reduce dependence on a single partition \citep{MeinshausenEtAl2009}. These references motivate the design but do not establish the cross-signed estimator in (\ref{Eqs: cross signed radius}); Proposition \ref{Prop: cross signed pilot radius} proves its rate. Because the $L_1$ functional is nondifferentiable at zero, the proposition makes no uniform local-to-zero claim \citep{HiranoPorter2012}.

Let
\[
\psi(C)=\frac{C-C_1}{C_m-C_1},
\]
and retain every raw pilot before projection so that values outside the candidate interval remain visible as endpoint diagnostics. For a sample size $N$, let $\check C_N$ be any pilot satisfying $\check C_N-C^\circ=O_{\mathbb P}(N^{-1/2})$, set $\check C_N^{\mathrm{proj}}=\Pi_{[C_1,C_m]}(\check C_N)$, and let $t_N=\psi(\check C_N^{\mathrm{proj}})$ and $t^\circ=\psi(C^\circ)$. For a deterministic sequence $\omega_N>0$, define the asymmetric radius discrepancy
\begin{equation}\label{Eqs: direct pilot discrepancy}
D_N(k)
=\{\psi(C_k)-t_N\}_+^2
+\omega_N\{t_N-\psi(C_k)\}_+^2,
\qquad [x]_+=\max(x,0).
\end{equation}

\begin{proposition}\label{Prop: asymmetric pilot target}
Suppose $t_N-t^\circ=O_{\mathbb P}(N^{-1/2})$, $\omega_N\to\infty$, and $\omega_N/N\to0$. Then
\[
\mathbb P\left[
\min\left\{\argmin_{k\in[m]}D_N(k)\right\}=k_C^*
\right]\longrightarrow1,
\]
where $k_C^*$ is defined in Proposition \ref{Prop: population pilot radius}.
\end{proposition}


\begin{proof}
Fix $k<k_C^*$. Strict increase of $\psi$ and the definition of $k_C^*$ give $\delta_k=t^\circ-\psi(C_k)>0$. Since $t_N\to_{\mathbb P}t^\circ$, with probability tending to one, $t_N-\psi(C_k)\ge\delta_k/2$. Hence $D_N(k)\ge\omega_N\delta_k^2/4\to\infty$ in probability.

For $k\ge k_C^*$, first suppose $C_{k_C^*}>C^\circ$. With probability tending to one, $t_N<\psi(C_{k_C^*})$, and therefore
\[
D_N(k)=\{\psi(C_k)-t_N\}^2
\xlongrightarrow{\mathbb P}
\{\psi(C_k)-t^\circ\}^2.
\]
The limiting value is uniquely minimised over $k\ge k_C^*$ at $k=k_C^*$ because $\psi$ is strictly increasing. If instead $C_{k_C^*}=C^\circ$, then
\[
D_N(k_C^*)
\le \max(1,\omega_N)(t_N-t^\circ)^2
=O_{\mathbb P}(\omega_N/N)=o_{\mathbb P}(1),
\]
whereas for every $k>k_C^*$,
\[
D_N(k)\xlongrightarrow{\mathbb P}
\{\psi(C_k)-t^\circ\}^2>0.
\]
Combining the underfitting and overfitting comparisons proves the result.
\end{proof}

The implementation uses
\[
\omega_N=
\left\{\log\!\log(N+\mathrm e^{\mathrm e})\right\}^{1/2},
\]
which is well defined and positive for every positive integer $N$, diverges as $N\to\infty$, and satisfies $\omega_N/N\to0$.

Proposition \ref{Prop: asymmetric pilot target} concerns direct minimisation of the discrepancy over the candidate radii. In the multiplier rule below, a radius is reachable only if it is selected by PanIC for at least one $\kappa\in\mathcal K_\kappa$; no assumption is made that every grid index is reachable. Accordingly, the proposition supplies the target-alignment rationale for the discrepancy, not a separate consistency theorem for the finite multiplier search. The first-order validity of the final full-sample selector follows instead from Proposition \ref{Prop: uniform kappa}.

\subsection{Repeated two-way cross-fitted multiplier rule}
Generate $R$ random balanced partitions $(\mathcal H_{r,1},\mathcal H_{r,2})$ of $[n]$, independently across $r=1,\ldots,R$, and use both directions of every partition. The split randomisation is independent of the observations, and $R$ remains fixed as $n\to\infty$. Index the resulting $B=2R$ calibration rows by $b$, with disjoint training and validation halves $(\mathcal T_b,\mathcal V_b)$ and sizes $n_{b,\mathrm{tr}}$ and $n_{b,\mathrm{val}}$. The implementation fixes $R=5$, giving ten rows, and both half sizes are asymptotic to $n/2$. If preprocessing is needed, its centring and scaling constants are estimated from $\mathcal T_b$ and applied to both halves. The rate statement then additionally requires the transformed ordinary estimators on both halves to satisfy the root-half-sample conditions in Proposition \ref{Prop: cross signed pilot radius} for the same fixed-coordinate $\beta^\circ$.

Let $\mathcal A_n$ be the measurable event that every ordinary half-sample empirical-risk minimiser used below exists uniquely and is finite, and that every constrained row fit returns finite coefficients and a finite empirical-risk value. The numerical implementation also treats a nonconverged or nonfinite solver result as $\mathcal A_n^c$. All ties and failure decisions use fixed deterministic rules.

Let $\mathcal K_\kappa=\{\kappa_1<\cdots<\kappa_L\}\subset[\underline\kappa,\overline\kappa]$ be fixed before inspecting the simulation outcomes, where $0<\underline\kappa\le\overline\kappa<\infty$. On $\mathcal A_n$, for each row compute ordinary unregularised fits on both halves,
\[
\widetilde\vartheta_b^{\mathrm{tr}}
=\argmin_{\vartheta\in\Theta_{\mathrm u}}R_{\mathcal T_b}(\vartheta),
\qquad
\widetilde\vartheta_b^{\mathrm{val}}
=\argmin_{\vartheta\in\Theta_{\mathrm u}}R_{\mathcal V_b}(\vartheta),
\]
and, for the Lasso calibration used here, define
\begin{align*}
\widetilde C_b^{\times}
&=\sum_{j=1}^d
\operatorname{sign}(\widetilde\beta_{b,j}^{\mathrm{tr}})
\widetilde\beta_{b,j}^{\mathrm{val}},\\
\widetilde C_b^{\mathrm{proj}}
&=\Pi_{[C_1,C_m]}(\widetilde C_b^{\times}),
\qquad
t_b=\psi(\widetilde C_b^{\mathrm{proj}}).
\end{align*}
The raw $\widetilde C_b^{\times}$ is retained for below-grid and above-grid diagnostics. On the training half, fit every constrained candidate once and record
\[
\widehat v_{b,k}=\min_{\vartheta\in\mathbb T_k}R_{\mathcal T_b}(\vartheta),
\qquad k\in[m].
\]
For each $\kappa\in\mathcal K_\kappa$, define
\[
\widehat k_b(\kappa)
=\min\left\{\argmin_{k\in[m]}
\{\widehat v_{b,k}+\kappa Q_{b,k}\}\right\},
\qquad
s_b(\kappa)=\psi\{C_{\widehat k_b(\kappa)}\},
\]
where $Q_{b,k}$ is the prescribed PanIC penalty shape computed at $n_{b,\mathrm{tr}}$. Assume $Q_{b,1}<\cdots<Q_{b,m}$ on the ordered candidate grid. The row score is
\begin{equation}\label{Eqs: fold radius discrepancy}
D_b(\kappa)
=\{s_b(\kappa)-t_b\}_+^2
+\omega_b\{t_b-s_b(\kappa)\}_+^2,
\qquad
\omega_b=
\left\{
\log\!\log(n_{b,\mathrm{val}}+\mathrm e^{\mathrm e})
\right\}^{1/2}.
\end{equation}
The larger weight on a selected radius below the pilot reflects the population ordering in Proposition \ref{Prop: population pilot radius}. The square-root log-log sequence is positive, diverges, and is $o(n_{b,\mathrm{val}})$. Relative to the log-log sequence, it moderates the finite-sample amplification of pilot error while preserving the two rate conditions used in Proposition \ref{Prop: asymmetric pilot target}.

For the rate interpretation, define measurable arbitrary completions of the half-sample estimators off $\mathcal A_n$ and use the same formulas to complete $\widetilde C_b^{\times}$, $\widetilde C_b^{\mathrm{proj}}$, and $t_b$. If $\mathbb P(\mathcal A_n)\to1$ and the completed estimators satisfy the two root-sample-size conditions of Proposition \ref{Prop: cross signed pilot radius} in every row, then fixed $B=2R$ and a finite-union argument give
\[
\max_{1\le b\le B}\left|\widetilde C_b^{\mathrm{proj}}-C^\circ\right|
=O_{\mathbb P}(n^{-1/2}).
\]
Because $\psi$ is affine, the same rate holds for $\max_b|t_b-t^\circ|$. Since $B$ is fixed and $n_{b,\mathrm{val}}\asymp n$,
\[
\max_{1\le b\le B}\omega_b
=O\!\left\{(\log\log n)^{1/2}\right\},
\]
and therefore
\[
\max_{1\le b\le B}\omega_b(t_b-t^\circ)^2
=O_{\mathbb P}\!\left\{
\frac{(\log\log n)^{1/2}}{n}
\right\}
=o_{\mathbb P}(1),
\]
as required by the asymmetric-discrepancy argument.

Define
\[
\overline D(\kappa)=\frac1B\sum_{b=1}^B D_b(\kappa),
\qquad
\widehat{\operatorname{se}}\{\overline D(\kappa)\}
=\left[
\frac{1}{B(B-1)}\sum_{b=1}^B
\{D_b(\kappa)-\overline D(\kappa)\}^2
\right]^{1/2}.
\]
Because rows from repeated partitions overlap and the two directions of a partition share observations, this is a descriptive across-row standard error for the one-standard-error rule, not an independent-row sampling standard error.
Let
\[
\kappa_{\min}
=\min\left\{\argmin_{\kappa\in\mathcal K_\kappa}\overline D(\kappa)\right\},
\]
and, on $\mathcal A_n$, select
\begin{equation}\label{Eqs: kappa one standard error}
\widehat\kappa_n^{\mathrm{cal}}
=\max\left\{\kappa\in\mathcal K_\kappa:
\overline D(\kappa)
\le \overline D(\kappa_{\min})
+\widehat{\operatorname{se}}\{\overline D(\kappa_{\min})\}
\right\}.
\end{equation}
Equation (\ref{Eqs: kappa one standard error}) defines $\widehat\kappa_n^{\mathrm{cal}}$ on $\mathcal A_n$. Fix $\kappa^\dagger\in\mathcal K_\kappa$ before observing the data (the implementation uses $\kappa^\dagger=1$) and define the total multiplier
\[
\widehat\kappa_n=
\begin{cases}
\widehat\kappa_n^{\mathrm{cal}},&\mathcal A_n,\\
\kappa^\dagger,&\mathcal A_n^c.
\end{cases}
\]
On $\mathcal A_n^c$ the procedure retains a calibration-failure flag; $\kappa^\dagger$ is a multiplier fallback and no regularised pilot is substituted. Thus $\widehat\kappa_n\in\mathcal K_\kappa\subset[\underline\kappa,\overline\kappa]$ on every outcome, so Proposition \ref{Prop: uniform kappa} applies directly. The cross-signed pilot-rate interpretation additionally requires $\mathbb P(\mathcal A_n)\to1$ and the half-sample rate conditions in Proposition \ref{Prop: cross signed pilot radius}; that success-probability condition is not needed for the uniform-multiplier conclusion because the fallback defines $\widehat\kappa_n$ on every outcome.
The particular admissible sequence used for $\omega_b$ can change which element of $\mathcal K_\kappa$ is selected, but it cannot affect this compact-range argument.
Finally, refit the constrained candidates on the full sample and select the smallest minimiser of their PanIC criteria with multiplier $\widehat\kappa_n$.

\begin{lemma}\label{Lemma: multiplier monotonicity}
Fix empirical values $u_1,\ldots,u_m$ and strictly increasing penalty shapes $Q_1<\cdots<Q_m$. Then
\[
k(\kappa)=\min\left\{\argmin_{k\in[m]}(u_k+\kappa Q_k)\right\}
\]
is nonincreasing in $\kappa>0$.
\end{lemma}

\begin{proof}
Let $0<\kappa_1<\kappa_2$ and write $k_j=k(\kappa_j)$. If $k_2>k_1$, optimality gives
\[
u_{k_1}+\kappa_1Q_{k_1}\le u_{k_2}+\kappa_1Q_{k_2},
\qquad
u_{k_2}+\kappa_2Q_{k_2}\le u_{k_1}+\kappa_2Q_{k_1}.
\]
Adding yields $(\kappa_2-\kappa_1)(Q_{k_2}-Q_{k_1})\le0$, contradicting strict increase of the $Q_k$. Thus $k_2\le k_1$.
\end{proof}

Lemma \ref{Lemma: multiplier monotonicity} gives the parsimony interpretation of the largest-multiplier rule in (\ref{Eqs: kappa one standard error}). A suggested multiplier grid is
\[
\mathcal K_{\kappa,0}
=\left\{10^{-2+4j/30}:j=0,\ldots,30\right\},
\]
although any alternative finite grid can also be used as long as it is fixed in advance, lies in a positive compact interval, and contains the prespecified default $\kappa^\dagger$. Frequent selection of either multiplier endpoint, raw cross-signed targets below $C_1$ or above $C_m$, or selected radii at either radius endpoint is an explicit diagnostic that the corresponding grid should be reconsidered.

For a full-sample fit using empirically estimated preprocessing, the fixed-model consistency conclusion additionally requires a justification for the resulting random parameter sets. The compact-multiplier argument alone does not supply it. The simulations avoid this issue by using the fixed predictor coordinates specified in Section \ref{Sect: Linear regression}.



\begin{remark}\label{Rem: pilot existence}
Calibration from the pilot scores requires a unique finite ordinary unregularised fit on both halves of every repeated split. In logistic regression, complete or quasi-complete separation can preclude a finite maximum-likelihood estimate; see \citet{albert_anderson_1984}. Analogous nonexistence or nonuniqueness may occur in other GLMs. If a required calibration fit fails, the procedure records a calibration failure and applies the prespecified $\kappa^\dagger$ convention to keep the multiplier defined. The numerical treatment of failed full-sample and cross-validation fits is separate, as described in Section \ref{Sect: Computations}.
\end{remark}
\end{bluechange}

\begin{bluechange}
\section{Reproducibility details}\label{App: reproducibility}
The simulations were implemented in R 4.4.0 using \texttt{glmnet} 4.1-10 and \texttt{Matrix} 1.7-0. The code, numerical settings, and retained outputs used for the tables and figures are available at \url{https://github.com/hiendn/PanIC_for_Regularised_GLMs}.
\end{bluechange}

\section{Proofs}


\begin{bluechange}
\subsection{Proof of Proposition \ref{Prop: uniform kappa}}
\begin{proof}
Let $E_n=\{\widehat\kappa_n\in[\underline\kappa,\overline\kappa]\}$. Since $\mathbb P(E_n)\to1$, for every $\eta>0$,
\[
\mathbb P\left(\max_kP_{k,n}>\eta\right)
\le \mathbb P(E_n^c)+
\mathbb P\left(\overline\kappa\max_kQ_{k,n}>\eta\right)\longrightarrow0.
\]
Thus B1 holds, and strict positivity follows from $\widehat\kappa_n>0$ and $Q_{k,n}>0$ almost surely. For $l>k$ and every $M>0$,
\begin{align*}
\mathbb P\{\sqrt n(P_{l,n}-P_{k,n})\le M\}
&\le \mathbb P(E_n^c)
 +\mathbb P\{\sqrt n(Q_{l,n}-Q_{k,n})\le M/\underline\kappa\}\longrightarrow0.
\end{align*}
This proves B2 without requiring the random difference $Q_{l,n}-Q_{k,n}$ to be nonnegative on every sample. The BIC-like statement is identical with $n$ in place of $\sqrt n$.
\end{proof}

\subsection{Proof of Proposition \ref{Prop: intercept envelopes}}
\begin{proof}
Along every segment joining two parameters in $\mathbb{U}\times\mathcal B_k$,
$|\eta|=|\beta_0+x^\top\beta|\le A_k(x)$. The gradients with respect to $\theta=(\beta_0,\beta)$ are
\begin{align*}
\nabla_\theta\ell_{\mathrm{lin}}&=-2(y-\eta)\widetilde x,\\
\nabla_\theta\ell_{\mathrm{logit}}&=\{h(\eta)-y\}\widetilde x,\\
\nabla_\theta\ell_{\mathrm{Pois}}&=\{e^\eta-y\}\widetilde x,\\
\nabla_\theta\ell_{\mathrm{Gamma}}&=\nu\{1-ye^{-\eta}\}\widetilde x.
\end{align*}
Here, in the logistic gradient, $h(\eta)=\{1+\exp(-\eta)\}^{-1}$.
The mean-value theorem, $|h(\eta)-y|\le1$ for $y\in\{0,1\}$, and $e^{\pm\eta}\le e^{A_k(x)}$ give the displayed envelopes. Fourth moments suffice in the linear case. For Poisson and Gamma regression, $\|X\|_2\leq\|X\|_1=\max_{s\in\{-1,1\}^d}s^\top X$, so $\exp(t\|X\|_2)\leq\sum_{s\in\{-1,1\}^d}\exp(ts^\top X)$. Every such exponential moment is finite for Gaussian or sub-Gaussian $X$; polynomial factors in $\|X\|_2$ are absorbed by increasing the exponential rate slightly. The Cauchy--Schwarz inequality then combines these bounds with $\mathbb E(Y^4)<\infty$. For A2, choose $(u_0,0)$ with $u_0\in\mathbb{U}$: the square of the baseline linear loss is controlled by $(|Y|+|u_0|)^4$, the logistic loss is bounded, the square of the Poisson loss is controlled using $\log(Y!)\le Y^2$ and the fourth moment of $Y$, and the square of the Gamma loss is controlled by a constant multiple of $1+Y^2+(\log Y)^2$.
\end{proof}

\subsection{Proof of Proposition \ref{Prop: estimated Gamma shape}}
\begin{proof}
Continuity and measurability follow on the fixed compact interval $\mathbb V=[\nu_-,\nu_+]$ with $0<\nu_-<\nu_+<\infty$. At a fixed baseline $(u_0,0,\nu_0)$, the loss in (\ref{Eqs: Gamma estimated loss}) is bounded by a constant multiple of $1+Y+|\log Y|$, so A2 holds. In addition to the $(\beta_0,\beta)$ gradient controlled in Proposition \ref{Prop: intercept envelopes},
\[
\frac{\partial\ell}{\partial\nu}
=\psi_0(\nu)-\log\nu-1+\eta-\log y+ye^{-\eta},
\]
where $\psi_0$ is the digamma function. The first three terms are uniformly bounded on $\mathbb{V}$, while $|\eta|\le A_k(x)$. Hence the full gradient is bounded by a constant multiple of
\[
1+|\log y|+A_k(x)+\{1+\|\widetilde x\|_2\}\{1+ye^{A_k(x)}\}.
\]
Its square is integrable by the sub-Gaussian exponential-moment bound, the Cauchy--Schwarz inequality, $\mathbb E(Y^4)<\infty$, and $\mathbb E\{(\log Y)^2\}<\infty$. Thus A1--A3 hold.
\end{proof}

\subsection{Proof of Proposition \ref{Prop: generic BIC penalty}}
\begin{proof}
For $n\ge2$, positivity follows from the definition; at $n=1$ it follows from the positive convention in the proposition. Since $m$ is fixed and $\psi(C_m)<\infty$,
\[
\max_kP_{k,n}\le\frac{\log n}{n}
\{c_0+\kappa_DO_{\mathbb P}(r_n)+\varepsilon\psi(C_m)\}
=o_{\mathbb P}(1).
\]
The cumulative maximum in (\ref{Eqs: cumulative complexity}) is nondecreasing. Thus, for $l>k$,
\begin{align*}
n(P_{l,n}-P_{k,n})
&=\log n\{\kappa_D(\widetilde D_{l,n}-\widetilde D_{k,n})
+\varepsilon(\psi(C_l)-\psi(C_k))\}\\
&\ge\varepsilon\log n\{\psi(C_l)-\psi(C_k)\}\longrightarrow\infty.
\end{align*}
This verifies B1 and $\mathrm{B2}^*$.
\end{proof}

\subsection{Proof of Theorem \ref{Thm: random finite grid}}
\begin{proof}
Fix $\epsilon>0$. If $c^*-\epsilon<a$, the underfitting event is empty. Otherwise, every $c\le c^*-\epsilon$ satisfies $v_n(c)\ge v_n(c^*-\epsilon)$. By nestedness, $v_n(c_n^+)\le v_n(c^*)$, while $P_n(c_n^+)\le P_n(b)=o_{\mathbb P}(1)$. Since $c^*$ is the smallest minimiser, $v(c^*-\epsilon)>v(c^*)$. The two convergences in (\ref{Eqs: random grid pointwise convergence}) therefore imply that, with probability tending to one, the criterion at every $c\le c^*-\epsilon$ exceeds that at $c_n^+$.

If $c^*+\epsilon>b$, the overfitting event is empty. Otherwise, on the event $c_n^+\le c^*+\epsilon/2$, every $c\ge c^*+\epsilon$ satisfies
\begin{align*}
&\{v_n(c)+P_n(c)\}-\{v_n(c_n^+)+P_n(c_n^+)\}\\
&\quad\ge v_n(b)-v_n(c^*)
+P_n(c^*+\epsilon)-P_n(c^*+\epsilon/2).
\end{align*}
The first term is $O_{\mathbb P}(a_n^{-1})$ by (\ref{Eqs: random grid empirical gap}); the second is positive and dominates $a_n^{-1}$ by (\ref{Eqs: random grid penalty separation}). Condition (\ref{Eqs: one sided coverage}) makes the event have probability tending to one. Both tails therefore vanish, proving $\widehat c_n^{\mathcal G}\to_{\mathbb P}c^*$.
\end{proof}

\subsection{Proof of Proposition \ref{Prop: scalar boundary}}
\begin{proof}
Set $Q_n=n^{-1}\sum_{i=1}^nX_i^2$ and $W_n=n^{-1/2}\sum_{i=1}^nX_i\epsilon_i$. The law of large numbers and the central limit theorem give $Q_n\to_{\mathbb P}q>0$ and $W_n\rightsquigarrow W\sim N(0,\tau^2)$. On $Q_n>0$, the ordinary least-squares estimate is $\widehat\beta_{\mathrm{OLS}}=\beta^\circ+W_n/(\sqrt n Q_n)$, and the constrained estimate is its projection onto $[-C,C]$. The event $Q_n=0$ has probability tending to zero and may use any fixed feasible selection. Consistency of the ordinary estimate makes the lower boundary irrelevant with probability tending to one. Consequently,
\[
\sqrt n(\widehat\beta_n-\beta^\circ)
=\min\{W_n/Q_n,0\}\rightsquigarrow\min\{W/q,0\}.
\]
Writing $U_n=\sqrt n(\widehat\beta_n-\beta^\circ)$, the exact quadratic identity is
\[
n\{R_n(\widehat\beta_n)-R_n(\beta^\circ)\}
=Q_nU_n^2-2W_nU_n
\rightsquigarrow-\frac{W^2}{q}\mathds 1\{W<0\}.
\]
Taking $Z=-2W$ and $h=2q$ gives the stated value limit. The critical cone is $(-\infty,0]$, so this also equals half the minimum of $2dZ+hd^2$ over $d\le0$.
\end{proof}

\subsection{Proof of Proposition \ref{Prop: subsampling fixed k}}
\begin{proof}
Fix $a$ and define the scalar estimator $T_n=a^\top\widehat\theta_{k,n}$ of $T=a^\top\theta_k^*$. Theorem \ref{Thm: Shapiro fixed k} gives
\[
\sqrt n(T_n-T)\rightsquigarrow a^\top\bar d_k(Z_k),
\]
whose distribution function is $F_{k,a}$. Apply \citet[Theorem 2.2.1]{politis1999subsampling} with normalising constants $\tau_n=\sqrt n$. The conditions $b\to\infty$ and $b/n\to0$ imply $\tau_b/\tau_n=\sqrt{b/n}\to0$, and the subsample statistic in that theorem is exactly $\sqrt{b_n}\,a^\top(\widehat\theta_{k,b_n,J}-\widehat\theta_{k,n})$. The theorem therefore yields $\widehat L_{n,b_n,a}(x)\to_{\mathbb P}F_{k,a}(x)$ at every continuity point of $F_{k,a}$. If only $B_n$ subsets are sampled, the stochastic-approximation result of \citet[Corollary 2.4.1]{politis1999subsampling} applies when $B_n\to\infty$ and gives the same limit.
\end{proof}
\end{bluechange}

\subsection{A bound for minimum empirical risks}
\label{App: rootn-mer}

\begin{lemma}
\label{Lem: rootn-mer}
Assume that $(X_i)_{i\in\mathbb{N}}$ is i.i.d. and that A1--A3 hold for some compact set $\mathbb{T}\subset\mathbb{R}^{q}$. Define
\[
m_n = \min_{\beta\in\mathbb{T}} R_n(\beta),\qquad m = \min_{\beta\in\mathbb{T}} r(\beta).
\]
Then
\[
\sqrt{n}\left(m_n-m\right)=O_{\mathbb{P}}(1).
\]
\end{lemma}

\begin{proof}
Under A1--A3, the class $\{\ell(\cdot;\beta):\beta\in\mathbb{T}\}$ satisfies the uniform central limit theorem used to obtain the limit of $\sqrt n(m_n-m)$. In particular, by \citet[Lemma 1]{nguyen_panic:_2023} (see also \citet[Fact 4.1]{pmlr-v235-westerhout24a}),
\[
\sqrt{n}\left(m_n-m\right)
\]
converges in distribution to a finite random variable. Boundedness in probability follows.
\end{proof}

\subsection{Proof of Theorem \ref{Thm: Infinite selection}}
\begin{proof}
Since $\mathcal{L}$ is a compact-valued continuous correspondence,
by the Berge Maximum Theorem \citep[Theorem 17.31]{aliprantis_infinite_2006},
the value function defined by
\[
m(k)=\min_{\beta\in\mathbb{T}_{k}}r(\beta)
\]
is continuous. A continuous function attains its minimum on a compact
set. Moreover, $\mathcal K$ is the inverse image of the attained minimum under the continuous function $m$, so it is a nonempty compact subset of $[a,b]$ and $k^*=\min\mathcal K$ is well-defined.

Next, fix $n\in\mathbb{N}$ and define the empirical value function
\[
m_{n}(k)=\min_{\beta\in\mathbb{T}_{k}}R_{n}(\beta).
\]
Under A1--A3 and the compact-valued continuity of $\mathcal{L}$,
the same Berge Maximum Theorem argument implies that $m_{n}$ is continuous
on $\left[a,b\right]$ (for each realisation of the data). Consider
the penalised objective
\[
F_{n}(k)=m_{n}(k)+P_{k,n}.
\]
Since both $m_n$ and $k\mapsto P_{k,n}(\omega)$ are continuous on $\left[a,b\right]$ for almost every realisation, the penalised objective $F_n$ is continuous and therefore attains its minimum on $\left[a,b\right]$. Hence, for each $n\in\mathbb{N}$, the argmin set in the definition
of $\hat{K}_{n}$ is non-empty and $\hat{K}_{n}$ is well-defined. For each fixed $k$, the measurable maximum theorem gives measurability of the criterion in the sample, while its sample paths are continuous in $k$; hence the criterion is jointly measurable by the Carath\'eodory measurability theorem. The argmin correspondence $\Gamma_n(\omega)=\argmin_{k\in[a,b]}F_n(k,\omega)$ is therefore nonempty, compact-valued, and measurable; see \citet[Theorem 18.19]{aliprantis_infinite_2006}. For every $t\in\mathbb R$,
\[
\{\min\Gamma_n<t\}=\{\Gamma_n\cap([a,b]\cap(-\infty,t))\ne\varnothing\},
\]
which is measurable by weak measurability of $\Gamma_n$. Hence the minimum tie-breaking selector $\widehat K_n=\min\Gamma_n$ is measurable.

To show convergence in probability, we will show for all $\sigma>0$,
\[
\mathbb{P}\left(\left|\hat{K}_{n}-k^{*}\right|\ge\sigma\right)\to0
\]
as $n\to\infty$. Fix $\sigma>0$.

First consider the case $\hat{K}_{n}\le k^{*}-\sigma$. If $\sigma>k^{*}-a$,
the statement is trivial. So assume $\sigma\le k^{*}-a$. Let $\varepsilon$
be
\[
\varepsilon=\frac{1}{2}\left\{ \min_{\beta\in\mathbb{T}_{k^{*}-\sigma}}r(\beta)-\min_{\beta\in\mathbb{T}_{k^{*}}}r(\beta)\right\} .
\]
Because $k^*$ is the smallest minimiser of the non-increasing value function $k\mapsto \min_{\beta\in\mathbb{T}_k} r(\beta)$, we have $\varepsilon>0$.
Conditions A1--A3 and the i.i.d. assumption are sufficient for the
uniform strong law of large numbers \citet[Thm. 9.60]{shapiro_lectures_2021},
which gives $R_{n}(\beta)\to r(\beta)$ uniformly on any compact set
$\mathbb{T}_{k}$ with probability one as $n\to\infty$. Using this
result, and assumption B1, for every $\delta>0$, there exists some
$N_{\delta}\in\mathbb{N}$ such that for all $n>N_{\delta}$, the
events
\begin{align*}
\left|\min_{\beta\in\mathbb{T}_{k^{*}-\sigma}}R_{n}(\beta)-\min_{\beta\in\mathbb{T}_{k^{*}-\sigma}}r(\beta)\right| & \le\frac{\varepsilon}{3},\\
\left|\min_{\beta\in\mathbb{T}_{k^{*}}}R_{n}(\beta)-\min_{\beta\in\mathbb{T}_{k^{*}}}r(\beta)\right| & \le\frac{\varepsilon}{3},
\end{align*}
and $P_{k^{*},n}\le\varepsilon/3$ occur with probability at least
$1-3\delta$. These events imply
\begin{align*}
\min_{\beta\in\mathbb{T}_{k^{*}-\sigma}}R_{n}(\beta) & \ge\min_{\beta\in\mathbb{T}_{k^{*}-\sigma}}r(\beta)-\frac{\varepsilon}{3}\\
 & =\min_{\beta\in\mathbb{T}_{k^{*}}}r(\beta)+\frac{5\varepsilon}{3}\\
 & \ge\min_{\beta\in\mathbb{T}_{k^{*}}}R_{n}(\beta)+P_{k^{*},n}+\varepsilon.
\end{align*}
So for all $k\in\left[a,k^{*}-\sigma\right]$,
\begin{align*}
\min_{\beta\in\mathbb{T}_{k^{*}}}R_{n}(\beta)+P_{k^{*},n} & <\min_{\beta\in\mathbb{T}_{k^{*}-\sigma}}R_{n}(\beta)+P_{a,n}\\
 & \le\min_{\beta\in\mathbb{T}_{k}}R_{n}(\beta)+P_{a,n}\\
 & \le\min_{\beta\in\mathbb{T}_{k}}R_{n}(\beta)+P_{k,n}
\end{align*}
occurs with probability at least $1-3\delta$. Taking $n\to\infty$
yields
\[
\mathbb{P}\left(\hat{K}_{n}\le k^{*}-\sigma\right)\to0.
\]

Now consider the case $\hat{K}_{n}\ge k^{*}+\sigma$. If $\sigma>b-k^{*}$,
the statement is trivial. So assume $\sigma\le b-k^{*}$. If B2 holds,
then Lemma \ref{Lem: rootn-mer} applied with $\mathbb{T}=\mathbb{T}_{b}$ yields
\[
\sqrt{n}\left(\min_{\beta\in\mathbb{T}_{b}}R_{n}(\beta)-\min_{\beta\in\mathbb{T}_{b}}r(\beta)\right)=O_{\mathbb{P}}(1),
\]
and similarly, Lemma \ref{Lem: rootn-mer} applied with $\mathbb{T}=\mathbb{T}_{k^{*}}$ yields
\[
\sqrt{n}\left(\min_{\beta\in\mathbb{T}_{k^{*}}}R_{n}(\beta)-\min_{\beta\in\mathbb{T}_{k^{*}}}r(\beta)\right)=O_{\mathbb{P}}(1).
\]
Since $\min_{\beta\in\mathbb{T}_{k^{*}}}r(\beta)=\min_{\beta\in\mathbb{T}_{b}}r(\beta)$,
it follows that
\[
\sqrt{n}\left(\min_{\beta\in\mathbb{T}_{b}}R_{n}(\beta)-\min_{\beta\in\mathbb{T}_{k^{*}}}R_{n}(\beta)\right)=O_{\mathbb{P}}(1).
\]
Fix $\delta>0$, there exists $M>0$ such that
\[
\sqrt{n}\left|\min_{\beta\in\mathbb{T}_{b}}R_{n}(\beta)-\min_{\beta\in\mathbb{T}_{k^{*}}}R_{n}(\beta)\right|\le M
\]
with probability at least $1-\delta$ for large $n$. By B2, for each
$\delta>0$ and $M>0$,
\[
\sqrt{n}\left\{ P_{k^{*}+\sigma,n}-P_{k^{*},n}\right\} >M
\]
with probability at least $1-\delta$ for large $n$. Therefore, with
probability at least $1-2\delta$,
\[
\min_{\beta\in\mathbb{T}_{k^{*}}}R_{n}(\beta)-\min_{\beta\in\mathbb{T}_{b}}R_{n}(\beta)\le\frac{M}{\sqrt{n}}<P_{k^{*}+\sigma,n}-P_{k^{*},n}.
\]
Thus for all $k\in\left[k^{*}+\sigma,b\right]$,
\begin{align*}
\min_{\beta\in\mathbb{T}_{k^{*}}}R_{n}(\beta)+P_{k^{*},n} & <\min_{\beta\in\mathbb{T}_{b}}R_{n}(\beta)+P_{k^{*}+\sigma,n}\\
 & \le\min_{\beta\in\mathbb{T}_{k}}R_{n}(\beta)+P_{k^{*}+\sigma,n}\\
 & \le\min_{\beta\in\mathbb{T}_{k}}R_{n}(\beta)+P_{k,n}.
\end{align*}
Hence $\lim_{n\to\infty}\mathbb{P}\left(\hat{K}_{n}-k^{*}\ge\sigma\right)=0$.

Alternatively, assume $\mathrm{B2}^{*}$ and C1--C5. For each $k\in\left[a,b\right]$,
let $\beta_{k}^{*}$ be the unique minimiser of $r$ over $\mathbb{T}_{k}$
from C1. \blue{Under A1--A3 and C1--C5, the cited minimum empirical-risk limit yields}
\[
n\left\{ \min_{\beta\in\mathbb{T}_{k}}R_{n}(\beta)-R_{n}(\beta_{k}^{*})\right\} =O_{\mathbb{P}}(1),
\]
see \citet[Lemma 3]{nguyen_panic:_2023}. Since $k^{*}$ is a minimiser
of $k\mapsto\min_{\beta\in\mathbb{T}_{k}}r(\beta)$ and the sets are
nested, we have
\[
r(\beta_{b}^{*})=\min_{\beta\in\mathbb{T}_{b}}r(\beta)=\min_{\beta\in\mathbb{T}_{k^{*}}}r(\beta)=r(\beta_{k^{*}}^{*}).
\]
Therefore C5 implies
\[
n\left\{ R_{n}(\beta_{b}^{*})-R_{n}(\beta_{k^{*}}^{*})\right\} =O_{\mathbb{P}}(1).
\]
Combining these displays gives
\[
n\left\{ \min_{\beta\in\mathbb{T}_{b}}R_{n}(\beta)-\min_{\beta\in\mathbb{T}_{k^{*}}}R_{n}(\beta)\right\} =O_{\mathbb{P}}(1).
\]
Fix $\delta>0$, there exists $M>0$ such that
\[
n\left|\min_{\beta\in\mathbb{T}_{b}}R_{n}(\beta)-\min_{\beta\in\mathbb{T}_{k^{*}}}R_{n}(\beta)\right|\le M
\]
with probability at least $1-\delta$ for large $n$. By $\mathrm{B2}^{*}$,
for each $\delta>0$ and $M>0$,
\[
n\left\{ P_{k^{*}+\sigma,n}-P_{k^{*},n}\right\} >M
\]
with probability at least $1-\delta$ for large $n$. Therefore, with
probability at least $1-2\delta$,
\[
\min_{\beta\in\mathbb{T}_{k^{*}}}R_{n}(\beta)-\min_{\beta\in\mathbb{T}_{b}}R_{n}(\beta)\le\frac{M}{n}<P_{k^{*}+\sigma,n}-P_{k^{*},n},
\]
and hence for all $k\in\left[k^{*}+\sigma,b\right]$,
\begin{align*}
\min_{\beta\in\mathbb{T}_{k^{*}}}R_{n}(\beta)+P_{k^{*},n} & <\min_{\beta\in\mathbb{T}_{b}}R_{n}(\beta)+P_{k^{*}+\sigma,n}\\
 & \le\min_{\beta\in\mathbb{T}_{k}}R_{n}(\beta)+P_{k^{*}+\sigma,n}\\
 & \le\min_{\beta\in\mathbb{T}_{k}}R_{n}(\beta)+P_{k,n}.
\end{align*}
Hence $\lim_{n\to\infty}\mathbb{P}\left(\hat{K}_{n}-k^{*}\ge\sigma\right)=0$.
The two tails together imply $\mathbb{P}(|\hat{K}_{n}-k^{*}|\ge\sigma)\to0$
for all $\sigma>0$, and hence (\ref{Eqs: convergence in probability})
holds.
\end{proof}

\subsection{Proof of Proposition \ref{Prop: Sufficient conditions}}
The componentwise calculations below are written in no-intercept Lasso notation for notational simplicity. Proposition \ref{Prop: intercept envelopes} supplies the explicit augmented-covariate bounds needed when the intercept is constrained to a compact interval, and the same compact-set argument applies to Ridge and Elastic-net slope constraints.
\subsubsection{Linear Regression}
\begin{proof}
    It suffices to verify the PanIC conditions,
    \begin{enumerate}
    \item[A1]: For each $(x,y)\in\mathbb R^d\times\mathbb R$, the map
    $\beta\mapsto(y-\beta^\top x)^2$ is continuous. For each fixed $\beta$, the same expression is Borel measurable as a function of $(x,y)$.
    \item[A2]: $\mathbb{T}_k = \{\beta\in \mathbb{R}^d: \lVert\beta\rVert_1 \leq C_k\}$ is compact. Set $\tau_k=0$. Since $\lVert\tau_k\rVert_1=0\leq C_k$, we have $\tau_k\in\mathbb{T}_k$ for every $k\in[m]$. This condition is satisfied if $Y$ has a finite fourth moment:
    \[\mathbb{E}\left(\ell_k(X, Y; \tau_k)^2\right) = \mathbb{E}\left(Y^4\right) < \infty.\]
    \item[A3]: Fix \blue{$(x,y)\in \mathbb R^d\times\mathbb R$}, $k \in \left[m\right]$ and $\beta, \beta' \in \mathbb{T}_k$. By the multivariate Mean Value Theorem \citep{davidson_stochastic_2021},
    \[|\ell_k(x, y; \beta) - \ell_k(x, y; \beta')| \leq \sup_{\beta^*} \left\lVert \left. \frac{\partial \ell_k}{\partial \beta}\right|_{\beta = \beta^*} \right \rVert_2 \lVert \beta - \beta'\rVert_2, \]
    where $\beta^*$ is any point on the line segment joining $\beta$ and $\beta'$. We can write $\beta^* = \lambda\beta + (1 - \lambda)\beta'$, for $\lambda\in\left[0, 1\right]$. Note that,
    \begin{align*}
        \left\lVert \left. \frac{\partial \ell_k}{\partial \beta}\right|_{\beta = \beta} \right \rVert_2 &= \lVert-2(y - \beta^\top x)x\rVert_2\\
        &= 2 \left\lvert y -  \beta^\top x\right\rvert\lVert x\rVert_2.
    \end{align*}
    This gives us
    \begin{align*}
        \sup_{\beta^*} \left\lVert \left. \frac{\partial \ell_k}{\partial \beta}\right|_{\beta = \beta^*} \right \rVert_2 &= \sup_{\lambda \in \left[0,1\right]} 2 \left\lvert y - (\lambda \beta  + (1 - \lambda)\beta')^\top x \right\rvert\lVert x\rVert_2.
    \end{align*}
    Using the triangle inequality, we obtain the bound
        \begin{align*}
        \sup_{\beta^*} \left\lVert \left. \frac{\partial \ell_k}{\partial \beta}\right|_{\beta = \beta^*} \right \rVert_2 &\leq \sup_{\lambda \in \left[0,1\right]} 2 \left(|y| + \lambda \left|\beta^\top x\right| + (1 - \lambda) \left|\beta'^\top x\right|\right)\lVert x\rVert_2\\
        &\leq \sup_{\lambda \in \left[0,1\right]} 2 \left(|y| + \lambda \left\lVert x\right\rVert_1 \lVert\beta\rVert_1 + (1 - \lambda) \left\lVert x\right\rVert_1 \lVert\beta'\rVert_1\right)\lVert x\rVert_2\\
        & \leq 2 \left(|y| + \left\lVert x\right\rVert_1 \lVert\beta\rVert_1 + \left\lVert x\right\rVert_1 \lVert\beta'\rVert_1\right)\lVert x\rVert_2\\
        & \leq 2 \left(|y| + \left\lVert x\right\rVert_1 \lVert\beta\rVert_1 + \left\lVert x\right\rVert_1 \lVert\beta'\rVert_1\right)\lVert x\rVert_1\\
        &\leq 2 \left(|y| + 2\left\lVert x\right\rVert_1 C_k\right)\lVert x\rVert_1.
    \end{align*}
    Let $\mathcal{G}_k(x, y) = 2 \left(|y| + 2\left\lVert x\right\rVert_1 C_k\right)\lVert x\rVert_1$, then the condition $\mathbb{E}\left(\mathcal{G}_k(X, Y)^2\right) < \infty$ is satisfied if both $X$ and $Y$ have a finite fourth moment.
\end{enumerate}
\end{proof}

\subsubsection{Logistic Regression}
\begin{proof}
    \begin{enumerate}
    \item[A1]: For each $(x,y)\in\mathbb R^d\times\{0,1\}$, the map
    \[
    \beta\mapsto \log\{1+\exp(\beta^\top x)\}-y\beta^\top x
    \]
    is continuous. For each fixed $\beta$, it is measurable as a function of $(x,y)$ on $\mathbb R^d\times\{0,1\}$ equipped with the product $\sigma$-algebra.
    \item[A2]: $\mathbb{T}_k$ is compact. Set $\tau_k=0$, which belongs to $\mathbb{T}_k$ for every $k\in[m]$. This condition is satisfied as
    \begin{align*}
        \mathbb{E}\left(\ell_k(X, Y; \tau_k)^2\right) &= \mathbb{E}\left(\left(-Y \log(1/2) - (1-Y)\log(1-1/2)\right)^2\right)\\
        &= \mathbb{E}\left(\left(-\log(1/2)\right)^2\right)\\
        &= (\log(1/2))^2 < \infty.
    \end{align*}
    \item[A3]: Fix \blue{$(x,y)\in \mathbb R^d\times\{0,1\}$}, $k \in \left[m\right]$ and $\beta, \beta'\in \mathbb{T}_k$. By the multivariate Mean Value Theorem, we have
    \[|\ell_k(x, y; \beta) - \ell_k(x, y; \beta')| \leq \sup_{\beta^*} \left\lVert \left. \frac{\partial \ell_k}{\partial \beta}\right|_{\beta = \beta^*} \right \rVert_2 \lVert \beta - \beta'\rVert_2, \]
    where $\beta^*$ is any point on the line segment joining $\beta$ and $\beta'$. We again write $\beta^* = \lambda\beta + (1 - \lambda)\beta'$, for $\lambda\in\left[0, 1\right]$. Note that,
    \begin{align*}
        \left.\frac{\partial \ell_k}{\partial \beta}\right|_{\beta = \beta} &= \left(\frac{1}{1 + e^{-\beta^\top x}} - y\right)x.
    \end{align*}
    Then,
    \begin{align*}
        \left\lVert \left. \frac{\partial \ell_k}{\partial \beta}\right|_{\beta = \beta} \right \rVert_2 &= \left\lvert \frac{1}{1 + e^{-\beta^\top x}} - y \right\rvert \left\lVert x\right\rVert_2\\
        &\leq \left\lVert x\right\rVert_2\\
        &\leq \left\lVert x\right\rVert_1.
    \end{align*}
    It follows that
    \begin{align*}
        \sup_{\beta^*} \left\lVert \left. \frac{\partial \ell_k}{\partial \beta}\right|_{\beta = \beta^*} \right \rVert_2 &\leq \left\lVert x\right\rVert_1.
    \end{align*}
    Let $\mathcal{G}_k(x, y) = \lVert x\rVert_1$, then the condition $\mathbb{E}\left(\mathcal{G}_k(X, Y)^2\right) < \infty$ is satisfied if $X$ has a finite second moment.
\end{enumerate}
\end{proof}

\subsubsection{Poisson Regression}
\begin{proof}
    \begin{enumerate}
    \item[A1]: For each $(x,y)\in\mathbb R^d\times\mathbb Z_{\ge0}$, the map
    \[
    \beta\mapsto \exp(\beta^\top x)-y\beta^\top x+\log(y!)
    \]
    is continuous. For each fixed $\beta$, it is measurable as a function of $(x,y)$ on $\mathbb R^d\times\mathbb Z_{\ge0}$ equipped with the product $\sigma$-algebra.
    \item[A2]: $\mathbb{T}_k$ is compact. Set $\tau_k=0$, which belongs to $\mathbb{T}_k$ for every $k\in[m]$. This condition is satisfied if $Y$ has a finite fourth moment. For any integer $y\in\mathbb{Z}_{\geq 0}$,
    \[
        \log(y!) = \sum_{j=1}^y \log(j) \leq \sum_{j=1}^y j = y(y+1)/2 \leq y^2.
    \]
    Hence,
    \begin{align*}
        \mathbb{E}\left(\ell_k(X, Y; \tau_k)^2\right) &= \mathbb{E}\left(\left(1+\log(Y!)\right)^2\right)\\
        &\leq \mathbb{E}\left(\left(1+ Y^2\right)^2\right)\\
        &= \mathbb{E}\left(Y^4 + 2Y^2 + 1\right) < \infty.
    \end{align*}
    \item[A3]: Fix \blue{$(x,y)\in \mathbb R^d\times\mathbb Z_{\ge0}$}, $k \in \left[m\right]$ and $\beta, \beta'\in \mathbb{T}_k$. By the multivariate Mean Value Theorem,
    \[|\ell_k(x, y; \beta) - \ell_k(x, y; \beta')| \leq \sup_{\beta^*} \left\lVert \left. \frac{\partial \ell_k}{\partial \beta}\right|_{\beta = \beta^*} \right \rVert_2 \lVert \beta - \beta'\rVert_2,\]
    where $\beta^*$ is any point on the line segment joining $\beta$ and $\beta'$. We write $\beta^* = \lambda\beta + (1 - \lambda)\beta'$, for $\lambda\in\left[0, 1\right]$. Note that,
    \begin{align*}
        \left.\frac{\partial \ell_k}{\partial \beta}\right|_{\beta=\beta} = \left(\exp( \beta^\top x) - y\right)x.
    \end{align*}
    Using triangle inequality,
    \begin{align*}
        \sup_{\beta^*} \left\lVert \left. \frac{\partial \ell_k}{\partial \beta}\right|_{\beta = \beta^*} \right \rVert_2 &= \sup_{\lambda\in\left[0, 1\right]}\left| \exp\left((\lambda\beta + (1 - \lambda)\beta')^\top x\right) - y\right| \lVert x\rVert_2\\
        &\leq \sup_{\lambda\in\left[0, 1\right]}\left( \exp\left(\left\lvert(\lambda\beta + (1 - \lambda)\beta')^\top x\right\rvert\right) + y\right) \lVert x\rVert_2\\
        &\leq \left( \exp\left(\lvert\beta^\top x\rvert + \lvert\beta'^\top x\rvert\right) + y\right) \lVert x\rVert_2\\
        &\leq \left( \exp(\lVert\beta\rVert_1\lVert x\rVert_1 + \lVert\beta'\rVert_1\lVert x\rVert_1) + y\right) \lVert x\rVert_2\\
        &\leq \left( \exp(2C_k\lVert x\rVert_1) + y\right) \lVert x\rVert_1.
    \end{align*}
    Let $\mathcal{G}_k(x, y) = \left( \exp(2C_k\lVert x\rVert_1) + y\right) \lVert x\rVert_1$, then the condition $\mathbb{E}\left(\mathcal{G}_k(X, Y)^2\right) < \infty$ is satisfied if
    \begin{align}\label{Eqs: Poisson condition}
        \mathbb{E}\left(\mathcal{G}_k(X, Y)^2\right) &= \mathbb{E}\left(\left( \exp(2C_k\lVert X\rVert_1) + Y\right)^2 \lVert X\rVert_1^2\right) < \infty.
    \end{align}
    To verify this under the stated assumptions, use the Cauchy--Schwarz inequality to obtain
    \begin{align*}
        \mathbb{E}\left(\left( \exp(2C_k\lVert X\rVert_1) + Y\right)^2 \lVert X\rVert_1^2\right)
        &\leq 2\mathbb{E}\left(\exp(4C_k\lVert X\rVert_1)\lVert X\rVert_1^2\right)\\
        &\quad + 2\left(\mathbb{E}(Y^4)\right)^{1/2}\left(\mathbb{E}\left(\exp(8C_k\lVert X\rVert_1)\lVert X\rVert_1^4\right)\right)^{1/2}.
    \end{align*}
    If $X$ is Gaussian or sub-Gaussian, then $\mathbb{E}\{\exp(t\lVert X\rVert_1)\}<\infty$ for every $t>0$, since
    \[
        \exp(t\lVert X\rVert_1)=\exp\!\left(t\max_{s\in\{-1,1\}^d}s^\top X\right)\leq \sum_{s\in\{-1,1\}^d}\exp(ts^\top X),
    \]
    and each linear functional $s^\top X$ is Gaussian or sub-Gaussian. Polynomial factors in $\lVert X\rVert_1$ are absorbed by increasing $t$ slightly, so the display is finite whenever $Y$ has a finite fourth moment.
\end{enumerate}
\end{proof}

\subsubsection{Gamma Regression}
\blue{For this fixed-shape calculation, we restore the constant $c_\nu=\log\Gamma(\nu)-\nu\log\nu$ omitted from Table \ref{Table: Regression}. Since $\nu$ is fixed, this changes neither the optimisers nor A1--A3.}
\begin{proof}
    \begin{enumerate}
    \item[A1]: For fixed $\nu>0$ and each $(x,y)\in\mathbb R^d\times(0,\infty)$, the map
    \[
    \beta\mapsto \nu\beta^\top x+\nu y\exp(-\beta^\top x)
    -(\nu-1)\log y+\log\Gamma(\nu)-\nu\log\nu
    \]
    is continuous. For each fixed $\beta$, it is Borel measurable as a function of $(x,y)$.
    \item[A2]: $\mathbb{T}_k$ is compact. Set $\tau_k=0$, which belongs to $\mathbb{T}_k$ for every $k\in[m]$. This condition is satisfied if $Y$ has a finite fourth moment and $\mathbb{E}\left\{\left(\log Y\right)^2\right\} < \infty$. Let $c_\nu = \log(\Gamma(\nu)) - \nu\log(\nu)$. Then
    \begin{align*}
        \mathbb{E}\left(\ell_k(X, Y; \tau_k)^2\right) &= \mathbb{E}\left(\left(c_\nu + \nu Y - (\nu - 1)\log Y \right)^2\right)\\
        &\leq 3c_\nu^2 + 3\nu^2\mathbb{E}\left(Y^2\right) + 3(\nu - 1)^2\mathbb{E}\left\{\left(\log Y\right)^2\right\}\\
        &< \infty.
    \end{align*}
    \item[A3]: Fix \blue{$(x,y)\in \mathbb R^d\times(0,\infty)$}, $k \in \left[m\right]$ and $\beta, \beta'\in \mathbb{T}_k$. By the multivariate Mean Value Theorem,
    \[|\ell_k(x, y; \beta) - \ell_k(x, y; \beta')| \leq \sup_{\beta^*} \left\lVert \left. \frac{\partial \ell_k}{\partial \beta}\right|_{\beta = \beta^*} \right \rVert_2 \lVert \beta - \beta'\rVert_2, \]
    where $\beta^*$ is any point on the line segment joining $\beta$ and $\beta'$. We write $\beta^* = \lambda\beta + (1 - \lambda)\beta'$, for $\lambda\in\left[0, 1\right]$.
    \begin{align*}
        \left.\frac{\partial \ell_k}{\partial \beta}\right|_{\beta=\beta} = \left(1 - y\exp(-\beta^\top x)\right)\nu x.
    \end{align*}
    Using triangle inequality,
    \begin{align*}
        \sup_{\beta^*} \left\lVert \left. \frac{\partial \ell_k}{\partial \beta}\right|_{\beta = \beta^*} \right \rVert_2 &= \sup_{\lambda\in\left[0, 1\right]} \left\lVert \left(1 - y\exp\left(-(\lambda\beta + (1 - \lambda)\beta')^\top x\right)\right)\nu x\right\rVert_2\\
        &= \sup_{\lambda\in\left[0, 1\right]} \nu\left\lvert1 - y\exp\left(-(\lambda\beta + (1 - \lambda)\beta')^\top x\right)\right\rvert \lVert x\rVert_2 \\
        &\leq \nu\left(1 + y\exp\left(\lvert\beta^\top x\rvert + \lvert\beta'^\top x\rvert\right)\right) \lVert x\rVert_1 \\
        &\leq \nu\left(1 + y\exp(\lVert\beta\rVert_1\lVert x\rVert_1 + \lVert\beta'\rVert_1\lVert x\rVert_1)\right) \lVert x\rVert_1\\
        &\leq \nu\left(1 + y\exp(2C_k\lVert x\rVert_1)\right) \lVert x\rVert_1.
    \end{align*}
    Let $\mathcal{G}_k(x, y) = \nu\left(1 + y\exp(2C_k\lVert x\rVert_1)\right) \lVert x\rVert_1$, we require the pair of covariates and response $(X, Y)$ to satisfy
    \begin{equation}\label{Eqs: Gamma condition}
    \mathbb{E}\left(\mathcal{G}_k(X, Y)^2\right) = \mathbb{E}\left(\nu^2 \left(1 + Y\exp(2C_k\lVert X\rVert_1)\right)^2 \lVert X\rVert_1^2 \right)< \infty.
    \end{equation}
    As in the Poisson case, the Cauchy--Schwarz inequality yields
    \begin{align*}
        \mathbb{E}\left(\nu^2 \left(1 + Y\exp(2C_k\lVert X\rVert_1)\right)^2 \lVert X\rVert_1^2 \right)
        &\leq 2\nu^2\mathbb{E}\left(\lVert X\rVert_1^2\right)\\
        &\quad + 2\nu^2\left(\mathbb{E}(Y^4)\right)^{1/2}\left(\mathbb{E}\left(\exp(8C_k\lVert X\rVert_1)\lVert X\rVert_1^4\right)\right)^{1/2}.
    \end{align*}
    If $X$ is Gaussian or sub-Gaussian, then the required exponential moments of $\lVert X\rVert_1$ are finite by the same argument as above. Hence (\ref{Eqs: Gamma condition}) holds whenever $Y$ has a finite fourth moment.
\end{enumerate}
\end{proof}


\subsection{Proof of Proposition \ref{Prop: BIC-like consistency}}

\begin{bluechange}
\begin{proof}
By Theorem 2 of \citet{nguyen_panic:_2023}, it is enough to verify C1--C5. We include the intercept. Write
\[
\widetilde X=(1,X^\top)^\top,
\qquad
\theta=(\beta_0,\beta^\top)^\top,
\qquad
\widetilde\Sigma=\mathbb E(\widetilde X\widetilde X^\top)
=\begin{bmatrix}1&\mu^\top\\ \mu&\Sigma+\mu\mu^\top\end{bmatrix},
\]
where $\mu=\mathbb E(X)$. For every $(a,b)\ne0$,
\[
\begin{bmatrix}a\\ b\end{bmatrix}^{\!\top}
\widetilde\Sigma
\begin{bmatrix}a\\ b\end{bmatrix}
=(a+b^\top\mu)^2+b^\top\Sigma b>0.
\]
Hence $\widetilde\Sigma$ is positive definite.

For candidate $k$, the population risk is
\[
r_k(\theta)=\mathbb E(Y^2)-2\theta^\top\mathbb E(\widetilde X Y)
+\theta^\top\widetilde\Sigma\theta,
\]
so its Hessian is $2\widetilde\Sigma$. The risk is therefore continuous and strictly convex. Each candidate parameter set is nonempty, compact, and convex, so it has a unique population minimiser. Proposition \ref{Prop: intercept envelopes} supplies an integrable Lipschitz modulus for the squared loss on that set. Integration of this modulus makes $r_k$ Lipschitz there. These facts verify C1. The squared loss is continuously differentiable in $\theta$, and the same envelope gives the required local Lipschitz control, which verifies C2.

For C3, the Lasso slope set is polyhedral. The Ridge constraint is twice continuously differentiable. When $C_k>0$, the zero slope is strictly feasible, so the Mangasarian--Fromovitz constraint qualification holds. When $C_k=0$, the slope set is the singleton $\{0\}$, which is second-order regular. For Elastic-net, use
\[
\|\beta\|_1=\max_{s\in\{-1,1\}^d}s^\top\beta.
\]
The constraint is equivalent to the finite system
\[
(1-\alpha)\|\beta\|_2^2+\alpha s^\top\beta\le C_k,
\qquad s\in\{-1,1\}^d.
\]
If $C_k>0$, the zero slope is strictly feasible. At any boundary point, the direction towards zero strictly decreases every active convex inequality. Thus the Mangasarian--Fromovitz condition holds for the active system, and second-order regularity follows from \citet[Proposition 3.88]{bonnans_perturbation_2000}. The case $C_k=0$ again gives the singleton $\{0\}$. Taking the Cartesian product with the common compact intercept interval preserves second-order regularity by \citet[Proposition 3.89]{bonnans_perturbation_2000}.

Positive definiteness of the Hessian gives quadratic growth at every candidate minimiser, which verifies C4. Finally, suppose two nested candidates have the same minimum population risk. The minimiser over the smaller set is feasible for the larger set and attains its minimum. Strict convexity makes the larger-set minimiser unique, so the two population minimisers coincide. Their empirical risks are then identical for every sample and every $n$, which verifies C5. The cited consistency theorem now applies.
\end{proof}
\end{bluechange}

\subsection{Proof of Lemma \ref{Lemma: continuity}}
\begin{proof}
Let
\[
\hat{\beta}=\arg\min_{\beta\in\mathbb{R}^d} f(\beta),
\]
which is well defined because the case $\lambda=0$ is included in the assumptions. Since $g$ is finite-valued and convex, it is continuous on $\mathbb{R}^d$; therefore the sublevel set
\[
\phi(\lambda)=\{\beta\in\mathbb{R}^d:g(\beta)\le g(\hat{\beta})\}=L_{g(\hat{\beta})}
\]
is closed and bounded, hence compact, for every $\lambda\in\mathbb{R}_{\ge 0}$. As a constant correspondence, $\phi$ is continuous.

Next, let
\[
L(\lambda,\beta)=f(\beta)+\lambda g(\beta),\qquad (\lambda,\beta)\in \mathbb{R}_{\ge 0}\times\mathbb{R}^d.
\]
The function $L$ is jointly continuous. For each $\lambda\ge 0$, let
\[
\beta_\lambda=\arg\min_{\beta\in\mathbb{R}^d} L(\lambda,\beta),
\]
which is unique by assumption. We first show that $\beta_\lambda\in\phi(\lambda)$ for all $\lambda\ge0$. If $\lambda=0$, then $\beta_0=\hat{\beta}$ by uniqueness. If $\lambda>0$, optimality of $\beta_\lambda$ gives
\[
f(\beta_\lambda)+\lambda g(\beta_\lambda)\le f(\hat{\beta})+\lambda g(\hat{\beta}),
\]
while optimality of $\hat{\beta}$ for $f$ gives $f(\hat{\beta})\le f(\beta_\lambda)$. Subtracting yields
\[
\lambda\{g(\beta_\lambda)-g(\hat{\beta})\}\le 0,
\]
and hence $g(\beta_\lambda)\le g(\hat{\beta})$.

Therefore $\beta_\lambda$ is also the unique minimiser of $L(\lambda,\cdot)$ over the compact feasible set $\phi(\lambda)$. By the Berge Maximum Theorem \citep[Theorem 17.31]{aliprantis_infinite_2006}, the argmin correspondence for the problem
\[
m(\lambda)=\min_{\beta\in\phi(\lambda)} L(\lambda,\beta)
\]
is upper hemicontinuous. Since it is single-valued, it is continuous as a function \citep[Lemma 17.6]{aliprantis_infinite_2006}. Thus $\lambda\mapsto \beta_\lambda$ is continuous. Finally, $\rho(\lambda)=g(\beta_\lambda)$ is continuous as the composition of the continuous maps $\lambda\mapsto\beta_\lambda$ and $g$.
\end{proof}

\subsection{Proof of Lemma \ref{Lemma: monotonicity}}
\begin{proof}
Let
\[
Q_\lambda(\beta)=f(\beta)+\lambda g(\beta).
\]
By optimality of $\beta_{\lambda_1}$ for $Q_{\lambda_1}$ and of $\beta_{\lambda_2}$ for $Q_{\lambda_2}$,
\begin{align}
f(\beta_{\lambda_1})+\lambda_1 g(\beta_{\lambda_1}) &\le f(\beta_{\lambda_2})+\lambda_1 g(\beta_{\lambda_2}), \label{Eqs: K1 inequality}\\
f(\beta_{\lambda_2})+\lambda_2 g(\beta_{\lambda_2}) &\le f(\beta_{\lambda_1})+\lambda_2 g(\beta_{\lambda_1}). \label{Eqs: K2 inequality}
\end{align}
Adding (\ref{Eqs: K1 inequality}) and (\ref{Eqs: K2 inequality}) gives
\[
(\lambda_2-\lambda_1)\{g(\beta_{\lambda_2})-g(\beta_{\lambda_1})\}\le 0.
\]
Since $\lambda_2>\lambda_1$, this implies
\[
g(\beta_{\lambda_2})\le g(\beta_{\lambda_1}).
\]

If the inequality is strict and $\lambda_1>0$, then (\ref{Eqs: K1 inequality}) yields
\[
f(\beta_{\lambda_1})-f(\beta_{\lambda_2})\le \lambda_1\{g(\beta_{\lambda_2})-g(\beta_{\lambda_1})\}<0,
\]
so $f(\beta_{\lambda_2})>f(\beta_{\lambda_1})$. If $\lambda_1=0$, strict inequality of the constraint values implies $\beta_{\lambda_2}\ne\beta_{\lambda_1}$; because $\beta_{\lambda_1}$ is the unique minimiser of $f$, the same strict conclusion follows.

Finally, if $g(\beta_{\lambda_2})=g(\beta_{\lambda_1})$, then both (\ref{Eqs: K1 inequality}) and (\ref{Eqs: K2 inequality}) force
\[
f(\beta_{\lambda_1})=f(\beta_{\lambda_2}),
\]
and therefore
\[
Q_{\lambda_1}(\beta_{\lambda_1})=Q_{\lambda_1}(\beta_{\lambda_2}).
\]
Because $Q_{\lambda_1}$ has a unique minimiser by assumption, we conclude that $\beta_{\lambda_2}=\beta_{\lambda_1}$.
\end{proof}

\subsection{Proof of Theorem \ref{Thm: Shapiro fixed k}}
\begin{proof}
Fix $k$, set $\theta=(\beta_0,\beta^\top)^\top$, $\theta_k^*=(\beta_{0,k}^*,\beta_k^{*\top})^\top$, and $\mathbb S=\mathbb T_k$. For squared loss,
\[
\nabla R_n(\theta_k^*)=-\frac{2}{n}\sum_{i=1}^n
\left(Y_i-\widetilde X_i^\top\theta_k^*\right)\widetilde X_i,
\qquad
\nabla^2R_n(\theta)=\frac{2}{n}\sum_{i=1}^n\widetilde X_i\widetilde X_i^\top,
\]
and the Hessian does not depend on $\theta$. \begin{bluechange}
The random vector formed by $Y^2$, $Y\widetilde X$, and the distinct entries of $\widetilde X\widetilde X^\top$ has finite covariance under the stated fourth-moment assumptions. Its multivariate central limit theorem, together with the quadratic representation of $R_n(\theta)$, implies convergence of $\sqrt n(R_n-r)$ in the supremum norm jointly with its gradient on the compact set $\mathbb S$. In particular, $\zeta_{k,n}\rightsquigarrow Z_k$. The law of large numbers also gives
\[
\nabla^2R_n\xlongrightarrow{\mathbb P}H=2\mathbb E(\widetilde X\widetilde X^\top).
\]
Because $H$ is positive definite, $r$ has quadratic growth at its unique constrained minimiser. Moreover, for every local increment $u$ the empirical objective has the exact quadratic expansion
\[
R_n(\theta_k^*+u)-R_n(\theta_k^*)
=u^\top\nabla R_n(\theta_k^*)+
\frac12u^\top\nabla^2R_n u.
\]
Together with second-order regularity of $\mathbb S$ and the interior condition on the intercept, these facts verify the hypotheses of \citet[Theorem 4.4]{shapiro_statistical_2000}. That theorem yields the value limits (\ref{Eqs: value limit})--(\ref{Eqs: value expansion}) with $\varphi_k$ defined in (\ref{Eqs: phi k}); uniqueness of the limiting minimiser then yields the solution limit (\ref{Eqs: solution limit}).
\end{bluechange}
\end{proof}

\bibliographystyle{plainnat}
\bibliography{citations}


\end{document}
